%%%%%%%%%%%%%%%%%%%%%%%%%%%%%%%%%%%%%%%%%
% Lachaise Assignment
% LaTeX Template
% Version 1.0 (26/6/2018)
%
% This template originates from:
% http://www.LaTeXTemplates.com
%
% Authors:
% Marion Lachaise & François Févotte
% Vel (vel@LaTeXTemplates.com)
%
% License:
% CC BY-NC-SA 3.0 (http://creativecommons.org/licenses/by-nc-sa/3.0/)
% 
%%%%%%%%%%%%%%%%%%%%%%%%%%%%%%%%%%%%%%%%%

%----------------------------------------------------------------------------------------
%	PACKAGES AND OTHER DOCUMENT CONFIGURATIONS
%----------------------------------------------------------------------------------------

\documentclass{article}

\usepackage[margin=1cm, bottom=2cm]{geometry}

% \input{structure.tex} % Include the file specifying the document structure and custom commands

% \usepackage{indentfirst}\setlength{\parindent}{0em}
\usepackage{listings}
\usepackage{float}
\usepackage{amsmath,amsfonts,stmaryrd,amssymb} % Math packages
\usepackage{graphicx}
\usepackage{subcaption}
\usepackage{tikz}

% add linking w/ blue color
\usepackage{xcolor}
\usepackage{hyperref}

\hypersetup{
    colorlinks=true,
    linkcolor=blue,
    filecolor=magenta,      
    urlcolor=blue,
    }
    

\numberwithin{equation}{section}
\numberwithin{figure}{section}

\title{ADI-FDTD for Inhomogeneous Media}

\begin{document}

\maketitle

\section{Preliminaries}

\subsection{Maxwell's Equations}
\begin{subequations}\label{eq:maxwell}
    \begin{align}
        \nabla \times \mathbf{H} & = \varepsilon \frac{\partial \mathbf{E}}{\partial t} + \sigma \mathbf{E}
        \label{eq:maxwell-ampere}                                                                           \\
        \nabla \times \mathbf{E} & = -\mu \frac{\partial \mathbf{H}}{\partial t}
        \label{eq:maxwell-faraday}
    \end{align}
\end{subequations}

Expanding \eqref{eq:maxwell-ampere} into components:
\begin{subequations}\label{eq:e-components}
    \begin{align}
        \frac{\partial H_z}{\partial y} - \frac{\partial H_y}{\partial z} & = \varepsilon \frac{\partial E_x}{\partial t} + \sigma E_x
        \label{eq:e-x-component}                                                                                                       \\
        \frac{\partial H_x}{\partial z} - \frac{\partial H_z}{\partial x} & = \varepsilon \frac{\partial E_y}{\partial t} + \sigma E_y
        \label{eq:e-y-component}                                                                                                       \\
        \frac{\partial H_y}{\partial x} - \frac{\partial H_x}{\partial y} & = \varepsilon \frac{\partial E_z}{\partial t} + \sigma E_z
        \label{eq:e-z-component}
    \end{align}
\end{subequations}

Expanding \eqref{eq:maxwell-faraday} into components:
\begin{subequations}\label{eq:h-components}
    \begin{align}
        \frac{\partial H_x}{\partial t} & = \frac{1}{\mu} \left( \frac{\partial E_y}{\partial z} - \frac{\partial E_z}{\partial y} \right)
        \label{eq:h-x-component}                                                                                                           \\
        \frac{\partial H_y}{\partial t} & = \frac{1}{\mu} \left( \frac{\partial E_z}{\partial x} - \frac{\partial E_x}{\partial z} \right)
        \label{eq:h-y-component}                                                                                                           \\
        \frac{\partial H_z}{\partial t} & = \frac{1}{\mu} \left( \frac{\partial E_x}{\partial y} - \frac{\partial E_y}{\partial x} \right)
        \label{eq:h-z-component}
    \end{align}
\end{subequations}

\subsection{Yee Grid in AMReX}

The Yee grid stores different field components at different physical locations within each computational cell. In AMReX this staggering is represented by the \texttt{IndexType} used to define each \texttt{MultiFab}: a cell-centered direction has index type 0 and is located at $i+\frac{1}{2}$, while a nodal direction has index type 1 and is located at $i$. The implementation uses
\begin{align*}
    \mathbf{E}_\alpha & : \text{cell-centered in the } \alpha \text{ direction and nodal in the other two directions}, \\
    \mathbf{H}_\alpha & : \text{nodal in the } \alpha \text{ direction and cell-centered in the other two directions}.
\end{align*}
Thus the component locations are
\begin{center}
    \begin{tabular}{r|lll}
        Component     & $x$ location    & $y$ location    & $z$ location    \\
        \hline
        $E_{x,i,j,k}$ & $i+\frac{1}{2}$ & $j$             & $k$             \\
        $E_{y,i,j,k}$ & $i$             & $j+\frac{1}{2}$ & $k$             \\
        $E_{z,i,j,k}$ & $i$             & $j$             & $k+\frac{1}{2}$ \\
        \hline
        $H_{x,i,j,k}$ & $i$             & $j+\frac{1}{2}$ & $k+\frac{1}{2}$ \\
        $H_{y,i,j,k}$ & $i+\frac{1}{2}$ & $j$             & $k+\frac{1}{2}$ \\
        $H_{z,i,j,k}$ & $i+\frac{1}{2}$ & $j+\frac{1}{2}$ & $k$
    \end{tabular}
\end{center}

The half-cell offsets therefore do not appear explicitly as array indices such as $i+\frac{1}{2}$. They are encoded in the \texttt{MultiFab} staggering. The integer offsets in the finite-difference stencil choose the two source values that bracket the field being updated. For example, $H_{x,i,j,k}$ is located at $j+\frac{1}{2}$ in the $y$ direction, while $E_{z,i,j,k}$ and $E_{z,i,j+1,k}$ are located at $j$ and $j+1$. Therefore
\begin{equation}
    \left.\frac{\partial E_z}{\partial y}\right|_{H_{x,i,j,k}} \approx \frac{E_{z,i,j+1,k}-E_{z,i,j,k}}{\Delta y}.
\end{equation}
Similarly, $E_{x,i,j,k}$ is located at $j$, while $H_{z,i,j-1,k}$ and $H_{z,i,j,k}$ are located at $j-\frac{1}{2}$ and $j+\frac{1}{2}$, so
\begin{equation}
    \left.\frac{\partial H_z}{\partial y}\right|_{E_{x,i,j,k}} \approx \frac{H_{z,i,j,k}-H_{z,i,j-1,k}}{\Delta y}.
\end{equation}
This is why the magnetic update below uses forward differences such as $j+1$ and $k+1$, while the electric update uses backward differences such as $j-1$ and $k-1$: each derivative is centered at the staggered location of the field component being updated.

\section{FDTD Method}

The explicit Yee FDTD update used here is written in three stages, matching the implementation: first update $\mathbf{H}$ by a half step, then update $\mathbf{E}$ by a full step, and finally update $\mathbf{H}$ by another half step. Define the update coefficients
\begin{equation}\label{eq:fdtd-coefficients}
    C_a = \frac{2\varepsilon - \sigma \Delta t}{2\varepsilon + \sigma \Delta t}, \qquad C_b = \frac{2\Delta t}{2\varepsilon + \sigma \Delta t}, \qquad D_b = \frac{\Delta t}{\mu},
\end{equation}
where $\Delta t$ is the time step, $\varepsilon$ is the permittivity, $\mu$ is the permeability, and $\sigma$ is the conductivity.

\subsection{First Magnetic Half-Step}
\begin{subequations}\label{eq:fdtd-first-h-half}
    \begin{align}
        H_{x,i,j+\frac{1}{2},k+\frac{1}{2}}^{n+\frac{1}{2}} & = H_{x,i,j+\frac{1}{2},k+\frac{1}{2}}^n - \frac{D_{b,i,j+\frac{1}{2},k+\frac{1}{2}}}{2} \left[\frac{E_{z,i,j+1,k+\frac{1}{2}}^n-E_{z,i,j,k+\frac{1}{2}}^n}{\Delta y} - \frac{E_{y,i,j+\frac{1}{2},k+1}^n-E_{y,i,j+\frac{1}{2},k}^n}{\Delta z}\right]
        \label{eq:fdtd-first-hx}                                                                                                                                                                                                                                                                                    \\
        H_{y,i+\frac{1}{2},j,k+\frac{1}{2}}^{n+\frac{1}{2}} & = H_{y,i+\frac{1}{2},j,k+\frac{1}{2}}^n - \frac{D_{b,i+\frac{1}{2},j,k+\frac{1}{2}}}{2} \left[\frac{E_{x,i+\frac{1}{2},j,k+1}^n-E_{x,i+\frac{1}{2},j,k}^n}{\Delta z} - \frac{E_{z,i+1,j,k+\frac{1}{2}}^n-E_{z,i,j,k+\frac{1}{2}}^n}{\Delta x}\right]
        \label{eq:fdtd-first-hy}                                                                                                                                                                                                                                                                                    \\
        H_{z,i+\frac{1}{2},j+\frac{1}{2},k}^{n+\frac{1}{2}} & = H_{z,i+\frac{1}{2},j+\frac{1}{2},k}^n - \frac{D_{b,i+\frac{1}{2},j+\frac{1}{2},k}}{2} \left[\frac{E_{y,i+1,j+\frac{1}{2},k}^n-E_{y,i,j+\frac{1}{2},k}^n}{\Delta x} - \frac{E_{x,i+\frac{1}{2},j+1,k}^n-E_{x,i+\frac{1}{2},j,k}^n}{\Delta y}\right].
        \label{eq:fdtd-first-hz}
    \end{align}
\end{subequations}

In the AMReX Yee grid, with vacuum media, the magnetic field updates are implemented as
\begin{subequations}\label{eq:fdtd-first-h-half-amrex}
    \begin{align}
        B_x(i, j, k) & = B_x(i, j, k) - \frac{1}{2\Delta t} \left[\frac{E_z(i, j + 1, k) - E_z(i, j, k)}{\Delta y} - \frac{E_y(i, j, k + 1) - E_y(i, j, k)}{\Delta z}\right]  \\
        B_y(i, j, k) & = B_y(i, j, k) - \frac{1}{2\Delta t} \left[\frac{E_x(i, j, k + 1) - E_x(i, j, k)}{\Delta z} - \frac{E_z(i + 1, j, k) - E_z(i, j, k)}{\Delta x}\right]  \\
        B_z(i, j, k) & = B_z(i, j, k) - \frac{1}{2\Delta t} \left[\frac{E_y(i + 1, j, k) - E_y(i, j, k)}{\Delta x} - \frac{E_x(i, j + 1, k) - E_x(i, j, k)}{\Delta y}\right].
    \end{align}
\end{subequations}

\subsection{Electric Full-Step}
\begin{subequations}\label{eq:fdtd-e-full}
    \begin{align}
        E_{x,i+\frac{1}{2},j,k}^{n+1} & = C_{a,i+\frac{1}{2},j,k}^{n+1} E_{x,i+\frac{1}{2},j,k}^n + C_{b,i+\frac{1}{2},j,k}^{n+1} \left[\frac{H_{z,i+\frac{1}{2},j+\frac{1}{2},k}^{n+\frac{1}{2}}-H_{z,i+\frac{1}{2},j-\frac{1}{2},k}^{n+\frac{1}{2}}}{\Delta y} - \frac{H_{y,i+\frac{1}{2},j,k+\frac{1}{2}}^{n+\frac{1}{2}}-H_{y,i+\frac{1}{2},j,k-\frac{1}{2}}^{n+\frac{1}{2}}}{\Delta z}\right]
        \label{eq:fdtd-ex}                                                                                                                                                                                                                                                                                                                                                                          \\
        E_{y,i,j+\frac{1}{2},k}^{n+1} & = C_{a,i,j+\frac{1}{2},k}^{n+1} E_{y,i,j+\frac{1}{2},k}^n + C_{b,i,j+\frac{1}{2},k}^{n+1} \left[\frac{H_{x,i,j+\frac{1}{2},k+\frac{1}{2}}^{n+\frac{1}{2}}-H_{x,i,j+\frac{1}{2},k-\frac{1}{2}}^{n+\frac{1}{2}}}{\Delta z} - \frac{H_{z,i+\frac{1}{2},j+\frac{1}{2},k}^{n+\frac{1}{2}}-H_{z,i-\frac{1}{2},j+\frac{1}{2},k}^{n+\frac{1}{2}}}{\Delta x}\right]
        \label{eq:fdtd-ey}                                                                                                                                                                                                                                                                                                                                                                          \\
        E_{z,i,j,k+\frac{1}{2}}^{n+1} & = C_{a,i,j,k+\frac{1}{2}}^{n+1} E_{z,i,j,k+\frac{1}{2}}^n + C_{b,i,j,k+\frac{1}{2}}^{n+1} \left[\frac{H_{y,i+\frac{1}{2},j,k+\frac{1}{2}}^{n+\frac{1}{2}}-H_{y,i-\frac{1}{2},j,k+\frac{1}{2}}^{n+\frac{1}{2}}}{\Delta x} - \frac{H_{x,i,j+\frac{1}{2},k+\frac{1}{2}}^{n+\frac{1}{2}}-H_{x,i,j-\frac{1}{2},k+\frac{1}{2}}^{n+\frac{1}{2}}}{\Delta y}\right].
        \label{eq:fdtd-ez}
    \end{align}
\end{subequations}

In the AMReX Yee grid, with vacuum media, the electric field updates are implemented as
\begin{subequations}\label{eq:fdtd-e-full-amrex}
    \begin{align}
        E_x(i,j,k) & = E_x(i,j,k) + c^2 \Delta t \left[\frac{B_z(i,j,k)-B_z(i,j-1,k)}{\Delta y} - \frac{B_y(i,j,k)-B_y(i,j,k-1)}{\Delta z}\right]  \\
        E_y(i,j,k) & = E_y(i,j,k) + c^2 \Delta t \left[\frac{B_x(i,j,k)-B_x(i,j,k-1)}{\Delta z} - \frac{B_z(i,j,k)-B_z(i-1,j,k)}{\Delta x}\right]  \\
        E_z(i,j,k) & = E_z(i,j,k) + c^2 \Delta t \left[\frac{B_y(i,j,k)-B_y(i-1,j,k)}{\Delta x} - \frac{B_x(i,j,k)-B_x(i,j-1,k)}{\Delta y}\right].
    \end{align}
\end{subequations}

\subsection{Second Magnetic Half-Step}
\begin{subequations}\label{eq:fdtd-second-h-half}
    \begin{align}
        H_{x,i,j+\frac{1}{2},k+\frac{1}{2}}^{n+1} & = H_{x,i,j+\frac{1}{2},k+\frac{1}{2}}^{n+\frac{1}{2}} - \frac{D_{b,i,j+\frac{1}{2},k+\frac{1}{2}}}{2} \left[\frac{E_{z,i,j+1,k+\frac{1}{2}}^{n+1}-E_{z,i,j,k+\frac{1}{2}}^{n+1}}{\Delta y} - \frac{E_{y,i,j+\frac{1}{2},k+1}^{n+1}-E_{y,i,j+\frac{1}{2},k}^{n+1}}{\Delta z}\right]
        \label{eq:fdtd-second-hx}                                                                                                                                                                                                                                                                                                       \\
        H_{y,i+\frac{1}{2},j,k+\frac{1}{2}}^{n+1} & = H_{y,i+\frac{1}{2},j,k+\frac{1}{2}}^{n+\frac{1}{2}} - \frac{D_{b,i+\frac{1}{2},j,k+\frac{1}{2}}}{2} \left[\frac{E_{x,i+\frac{1}{2},j,k+1}^{n+1}-E_{x,i+\frac{1}{2},j,k}^{n+1}}{\Delta z} - \frac{E_{z,i+1,j,k+\frac{1}{2}}^{n+1}-E_{z,i,j,k+\frac{1}{2}}^{n+1}}{\Delta x}\right]
        \label{eq:fdtd-second-hy}                                                                                                                                                                                                                                                                                                       \\
        H_{z,i+\frac{1}{2},j+\frac{1}{2},k}^{n+1} & = H_{z,i+\frac{1}{2},j+\frac{1}{2},k}^{n+\frac{1}{2}} - \frac{D_{b,i+\frac{1}{2},j+\frac{1}{2},k}}{2} \left[\frac{E_{y,i+1,j+\frac{1}{2},k}^{n+1}-E_{y,i,j+\frac{1}{2},k}^{n+1}}{\Delta x} - \frac{E_{x,i+\frac{1}{2},j+1,k}^{n+1}-E_{x,i+\frac{1}{2},j,k}^{n+1}}{\Delta y}\right].
        \label{eq:fdtd-second-hz}
    \end{align}
\end{subequations}

In the AMReX Yee grid, with vacuum media, the second magnetic field updates are implemented as
\begin{subequations}\label{eq:fdtd-second-h-half-amrex}
    \begin{align}
        B_x(i,j,k) & = B_x(i,j,k) - \frac{1}{2\Delta t} \left[\frac{E_z(i,j+1,k)-E_z(i,j,k)}{\Delta y} - \frac{E_y(i,j,k+1)-E_y(i,j,k)}{\Delta z}\right]  \\
        B_y(i,j,k) & = B_y(i,j,k) - \frac{1}{2\Delta t} \left[\frac{E_x(i,j,k+1)-E_x(i,j,k)}{\Delta z} - \frac{E_z(i+1,j,k)-E_z(i,j,k)}{\Delta x}\right]  \\
        B_z(i,j,k) & = B_z(i,j,k) - \frac{1}{2\Delta t} \left[\frac{E_y(i+1,j,k)-E_y(i,j,k)}{\Delta x} - \frac{E_x(i,j+1,k)-E_x(i,j,k)}{\Delta y}\right].
    \end{align}
\end{subequations}

This explicit method is simple and local, but the time step is restricted by the CFL stability condition. For a uniform lossless grid, a typical stability limit is
\begin{equation}\label{eq:fdtd-cfl}
    \Delta t \le \frac{1}{c\sqrt{\frac{1}{\Delta x^2}+\frac{1}{\Delta y^2}+\frac{1}{\Delta z^2}}}, \qquad c = \frac{1}{\sqrt{\mu\varepsilon}}.
\end{equation}


\section{ADI Method}

The ADI-FDTD update advances the fields by a full time step $\Delta t$ in two half steps of size $\Delta t/2$. Unlike the explicit three-stage leapfrog update above, each ADI half step first updates $\mathbf{E}$ implicitly over $\Delta t/2$, then updates $\mathbf{H}$ explicitly over $\Delta t/2$ from the updated electric field. The first half step treats one set of implicit directions; the second half step treats the complementary directions to complete the cycle. Define the half-step update coefficients
\begin{equation}\label{eq:adi-half-coefficients}
    C_a = \frac{4\varepsilon-\sigma\Delta t}{4\varepsilon+\sigma\Delta t}, \qquad C_b = \frac{2\Delta t}{4\varepsilon+\sigma\Delta t}, \qquad D_b = \frac{\Delta t}{2\mu},
\end{equation}
where $\Delta t$ is the time step, $\varepsilon$ is the permittivity, $\mu$ is the permeability, and $\sigma$ is the conductivity. These differ from the full-step coefficients in \eqref{eq:fdtd-coefficients} because each electric update uses $\Delta t/2$.

\subsection{First ADI Half-Step}

\paragraph{$E_x$ and $H_z$}
By equation \eqref{eq:e-x-component},
\begin{equation}
    \varepsilon_{i+\tfrac12,j,k} \frac{E_{x,i+\tfrac12,j,k}^{n+\frac12}-E_{x,i+\tfrac12,j,k}^{n}}{\Delta t/2} + \sigma_{i+\tfrac12,j,k} \frac{E_{x,i+\tfrac12,j,k}^{n+\frac12}+E_{x,i+\tfrac12,j,k}^{n}}{2} = [H_1],
    \label{eq:first-e-ex}
\end{equation}
where
\begin{equation}
    [H_1] = \frac{H_{z,i+\tfrac12,j+\tfrac12,k}^{n+\frac12}-H_{z,i+\tfrac12,j-\tfrac12,k}^{n+\frac12}}{\Delta y} - \frac{H_{y,i+\tfrac12,j,k+\tfrac12}^{n}-H_{y,i+\tfrac12,j,k-\tfrac12}^{n}}{\Delta z}.
\end{equation}
Rearranging gives
\begin{equation}
    \left(\frac{2\varepsilon}{\Delta t}+\frac{\sigma}{2}\right) E_{x}^{n+\frac12} = \left(\frac{2\varepsilon}{\Delta t}-\frac{\sigma}{2}\right) E_{x}^{n} + [H_1].
\end{equation}
Hence
\begin{equation}
    E_{x,i+\tfrac12,j,k}^{n+\frac12} = C_{a,i+\tfrac12,j,k}\,E_{x,i+\tfrac12,j,k}^{n} + C_{b,i+\tfrac12,j,k}\,[H_1],
    \label{eq:first-e-ex-adi}
\end{equation}

By equation \eqref{eq:h-z-component},
\begin{align}
    \begin{split}
        H_{z,i+\tfrac12,j+\tfrac12,k}^{n+\frac12} = H_{z,i+\tfrac12,j+\tfrac12,k}^{n} + D_{b,i+\tfrac12,j+\tfrac12,k}\Big( & \frac{E_{x,i+\tfrac12,j+1,k}^{n+\frac12}-E_{x,i+\tfrac12,j,k}^{n+\frac12}}{\Delta y} \\
                                                                                                                           & - \frac{E_{y,i+1,j+\tfrac12,k}^{n}-E_{y,i,j+\tfrac12,k}^{n}}{\Delta x}\Big),
    \end{split}
\end{align}

Substituting the magnetic update into \eqref{eq:first-e-ex-adi} and eliminating $H_{z}^{n+\frac12}$ terms gives
\begin{align}
    \begin{split}
        E_{x,i+\tfrac12,j,k}^{n+\frac12} & = C_{a,i+\tfrac12,j,k}\,E_{x,i+\tfrac12,j,k}^{n}                                                                                                                                                                                                                                                           \\
                                         & \quad + \frac{C_{b,i+\tfrac12,j,k}}{\Delta y}\Bigg[H_{z,i+\tfrac12,j+\tfrac12,k}^{n} + D_{b,i+\tfrac12,j+\tfrac12,k}\left(\frac{E_{x,i+\tfrac12,j+1,k}^{n+\frac12}-E_{x,i+\tfrac12,j,k}^{n+\frac12}}{\Delta y}\right. \left. - \frac{E_{y,i+1,j+\tfrac12,k}^{n}-E_{y,i,j+\tfrac12,k}^{n}}{\Delta x}\right) \\
                                         & \qquad\qquad\qquad - H_{z,i+\tfrac12,j-\tfrac12,k}^{n} - D_{b,i+\tfrac12,j-\tfrac12,k}\left(\frac{E_{x,i+\tfrac12,j,k}^{n+\frac12}-E_{x,i+\tfrac12,j-1,k}^{n+\frac12}}{\Delta y}\right. \left. - \frac{E_{y,i+1,j-\tfrac12,k}^{n}-E_{y,i,j-\tfrac12,k}^{n}}{\Delta x}\right)\Bigg]                         \\
                                         & \quad - \frac{C_{b,i+\tfrac12,j,k}}{\Delta z}\left(H_{y,i+\tfrac12,j,k+\tfrac12}^{n}-H_{y,i+\tfrac12,j,k-\tfrac12}^{n}\right).
    \end{split}
\end{align}

Collecting the unknown $E_{x}^{n+\frac12}$ terms gives
\begin{equation}
    \begin{aligned}
         & \quad -\frac{C_{b,i+\tfrac12,j,k}\,D_{b,i+\tfrac12,j+\tfrac12,k}}{\Delta y^2}\, E_{x,i+\tfrac12,j+1,k}^{n+\frac12}                                                                                     \\
         & \quad + \left[1+\frac{C_{b,i+\tfrac12,j,k}\,D_{b,i+\tfrac12,j+\tfrac12,k}}{\Delta y^2} +\frac{C_{b,i+\tfrac12,j,k}\,D_{b,i+\tfrac12,j-\tfrac12,k}}{\Delta y^2}\right] E_{x,i+\tfrac12,j,k}^{n+\frac12} \\
         & \quad - \frac{C_{b,i+\tfrac12,j,k}\,D_{b,i+\tfrac12,j-\tfrac12,k}}{\Delta y^2}\, E_{x,i+\tfrac12,j-1,k}^{n+\frac12}                                                                                    \\
         & = C_{a,i+\tfrac12,j,k}\, E_{x,i+\tfrac12,j,k}^{n} + \frac{C_{b,i+\tfrac12,j,k}}{\Delta y} \Big(H_{z,i+\tfrac12,j+\tfrac12,k}^{n}-H_{z,i+\tfrac12,j-\tfrac12,k}^{n}\Big)                                \\
         & \quad - \frac{C_{b,i+\tfrac12,j,k}\,D_{b,i+\tfrac12,j+\tfrac12,k}}{\Delta y\,\Delta x} \Big(E_{y,i+1,j+\tfrac12,k}^{n}-E_{y,i,j+\tfrac12,k}^{n}\Big)                                                   \\
         & \quad + \frac{C_{b,i+\tfrac12,j,k}\,D_{b,i+\tfrac12,j-\tfrac12,k}}{\Delta y\,\Delta x} \Big(E_{y,i+1,j-\tfrac12,k}^{n}-E_{y,i,j-\tfrac12,k}^{n}\Big)                                                   \\
         & \quad - \frac{C_{b,i+\tfrac12,j,k}}{\Delta z} \Big(H_{y,i+\tfrac12,j,k+\tfrac12}^{n}-H_{y,i+\tfrac12,j,k-\tfrac12}^{n}\Big).
    \end{aligned}
\end{equation}

Equivalently, the tridiagonal system can be written as
\begin{align}
    \begin{split}
        -\alpha_{x1}E_{x,i+\tfrac12,j-1,k}^{n+\frac12} + \beta_{x1}E_{x,i+\tfrac12,j,k}^{n+\frac12} - \gamma_{x1}E_{x,i+\tfrac12,j+1,k}^{n+\frac12}
         & = p_{x1}E_{x,i+\tfrac12,j,k}^{n}                                                                       \\
         & \quad +\frac{1}{\Delta y}\Big(H_{z,i+\tfrac12,j+\tfrac12,k}^{n}-H_{z,i+\tfrac12,j-\tfrac12,k}^{n}\Big) \\
         & \quad -\frac{1}{\Delta z}\Big(H_{y,i+\tfrac12,j,k+\tfrac12}^{n}-H_{y,i+\tfrac12,j,k-\tfrac12}^{n}\Big) \\
         & \quad +\frac{q_{x1}}{\Delta x\Delta y}\Big(E_{y,i+1,j-\tfrac12,k}^{n}-E_{y,i,j-\tfrac12,k}^{n}\Big)    \\
         & \quad -\frac{r_{x1}}{\Delta x\Delta y}\Big(E_{y,i+1,j+\tfrac12,k}^{n}-E_{y,i,j+\tfrac12,k}^{n}\Big),
    \end{split}
    \label{eq:first-e-ex-adi-tridiagonal}
\end{align}

and the magnetic update can be written as
\begin{equation}
    H_{z,i+\tfrac12,j+\tfrac12,k}^{n+\frac12} = H_{z,i+\tfrac12,j+\tfrac12,k}^{n} + \frac{r_{x1}}{\Delta y}\Big(E_{x,i+\tfrac12,j+1,k}^{n+\frac12}-E_{x,i+\tfrac12,j,k}^{n+\frac12}\Big) - \frac{r_{x1}}{\Delta x}\Big(E_{y,i+1,j+\tfrac12,k}^{n}-E_{y,i,j+\tfrac12,k}^{n}\Big).
    \label{eq:adi-first-hz}
\end{equation}
where
\begin{align}
    \alpha_{x1} & =\frac{D_{b,i+\tfrac12,j-\tfrac12,k}}{\Delta y^2},       \\
    \beta_{x1}  & =\frac{1}{C_{b,i+\tfrac12,j,k}}+\alpha_{x1}+\gamma_{x1}, \\
    \gamma_{x1} & =\frac{D_{b,i+\tfrac12,j+\tfrac12,k}}{\Delta y^2},       \\
    p_{x1}      & =\frac{C_{a,i+\tfrac12,j,k}}{C_{b,i+\tfrac12,j,k}},      \\
    q_{x1}      & =D_{b,i+\tfrac12,j-\tfrac12,k},                          \\
    r_{x1}      & =D_{b,i+\tfrac12,j+\tfrac12,k}.
\end{align}

\paragraph{$E_y$ and $H_x$}
By equation \eqref{eq:e-y-component},
\begin{equation}
    \varepsilon_{i,j+\tfrac12,k} \frac{E_{y,i,j+\tfrac12,k}^{n+\frac12}-E_{y,i,j+\tfrac12,k}^{n}}{\Delta t/2} + \sigma_{i,j+\tfrac12,k} \frac{E_{y,i,j+\tfrac12,k}^{n+\frac12}+E_{y,i,j+\tfrac12,k}^{n}}{2} = [H_2],
    \label{eq:first-e-ey}
\end{equation}
where
\begin{equation}
    [H_2] = \frac{H_{x,i,j+\tfrac12,k+\tfrac12}^{n+\frac12}-H_{x,i,j+\tfrac12,k-\tfrac12}^{n+\frac12}}{\Delta z} - \frac{H_{z,i+\tfrac12,j+\tfrac12,k}^{n}-H_{z,i-\tfrac12,j+\tfrac12,k}^{n}}{\Delta x}.
\end{equation}
Rearranging gives
\begin{equation}
    \left(\frac{2\varepsilon}{\Delta t}+\frac{\sigma}{2}\right) E_{y}^{n+\frac12} = \left(\frac{2\varepsilon}{\Delta t}-\frac{\sigma}{2}\right) E_{y}^{n} + [H_2].
\end{equation}
Hence
\begin{equation}
    E_{y,i,j+\tfrac12,k}^{n+\frac12} = C_{a,i,j+\tfrac12,k}\,E_{y,i,j+\tfrac12,k}^{n} + C_{b,i,j+\tfrac12,k}\,[H_2],
    \label{eq:first-e-ey-adi}
\end{equation}

By equation \eqref{eq:h-x-component},
\begin{align}
    \begin{split}
        H_{x,i,j+\tfrac12,k+\tfrac12}^{n+\frac12} = H_{x,i,j+\tfrac12,k+\tfrac12}^{n} + D_{b,i,j+\tfrac12,k+\tfrac12}\Big( & \frac{E_{y,i,j+\tfrac12,k+1}^{n+\frac12}-E_{y,i,j+\tfrac12,k}^{n+\frac12}}{\Delta z} \\
                                                                                                                           & - \frac{E_{z,i,j+1,k+\tfrac12}^{n}-E_{z,i,j,k+\tfrac12}^{n}}{\Delta y}\Big),
    \end{split}
\end{align}

Substituting the magnetic update into \eqref{eq:first-e-ey-adi} and eliminating $H_{x}^{n+\frac12}$ terms gives
\begin{align}
    \begin{split}
        E_{y,i,j+\tfrac12,k}^{n+\frac12} & = C_{a,i,j+\tfrac12,k}\,E_{y,i,j+\tfrac12,k}^{n}                                                                                                                                                                                                                                                           \\
                                         & \quad + \frac{C_{b,i,j+\tfrac12,k}}{\Delta z}\Bigg[H_{x,i,j+\tfrac12,k+\tfrac12}^{n} + D_{b,i,j+\tfrac12,k+\tfrac12}\left(\frac{E_{y,i,j+\tfrac12,k+1}^{n+\frac12}-E_{y,i,j+\tfrac12,k}^{n+\frac12}}{\Delta z}\right. \left. - \frac{E_{z,i,j+1,k+\tfrac12}^{n}-E_{z,i,j,k+\tfrac12}^{n}}{\Delta y}\right) \\
                                         & \qquad\qquad\qquad - H_{x,i,j+\tfrac12,k-\tfrac12}^{n} - D_{b,i,j+\tfrac12,k-\tfrac12}\left(\frac{E_{y,i,j+\tfrac12,k}^{n+\frac12}-E_{y,i,j+\tfrac12,k-1}^{n+\frac12}}{\Delta z}\right. \left. - \frac{E_{z,i,j+1,k-\tfrac12}^{n}-E_{z,i,j,k-\tfrac12}^{n}}{\Delta y}\right)\Bigg]                         \\
                                         & \quad - \frac{C_{b,i,j+\tfrac12,k}}{\Delta x}\left(H_{z,i+\tfrac12,j+\tfrac12,k}^{n}-H_{z,i-\tfrac12,j+\tfrac12,k}^{n}\right).
    \end{split}
\end{align}

Collecting the unknown $E_{y}^{n+\frac12}$ terms gives
\begin{equation}
    \begin{aligned}
         & \quad -\frac{C_{b,i,j+\tfrac12,k}\,D_{b,i,j+\tfrac12,k+\tfrac12}}{\Delta z^2}\, E_{y,i,j+\tfrac12,k+1}^{n+\frac12}                                                                                     \\
         & \quad + \left[1+\frac{C_{b,i,j+\tfrac12,k}\,D_{b,i,j+\tfrac12,k+\tfrac12}}{\Delta z^2} +\frac{C_{b,i,j+\tfrac12,k}\,D_{b,i,j+\tfrac12,k-\tfrac12}}{\Delta z^2}\right] E_{y,i,j+\tfrac12,k}^{n+\frac12} \\
         & \quad - \frac{C_{b,i,j+\tfrac12,k}\,D_{b,i,j+\tfrac12,k-\tfrac12}}{\Delta z^2}\, E_{y,i,j+\tfrac12,k-1}^{n+\frac12}                                                                                    \\
         & = C_{a,i,j+\tfrac12,k}\, E_{y,i,j+\tfrac12,k}^{n} + \frac{C_{b,i,j+\tfrac12,k}}{\Delta z} \Big(H_{x,i,j+\tfrac12,k+\tfrac12}^{n}-H_{x,i,j+\tfrac12,k-\tfrac12}^{n}\Big)                                \\
         & \quad - \frac{C_{b,i,j+\tfrac12,k}\,D_{b,i,j+\tfrac12,k+\tfrac12}}{\Delta z\,\Delta y} \Big(E_{z,i,j+1,k+\tfrac12}^{n}-E_{z,i,j,k+\tfrac12}^{n}\Big)                                                   \\
         & \quad + \frac{C_{b,i,j+\tfrac12,k}\,D_{b,i,j+\tfrac12,k-\tfrac12}}{\Delta z\,\Delta y} \Big(E_{z,i,j+1,k-\tfrac12}^{n}-E_{z,i,j,k-\tfrac12}^{n}\Big)                                                   \\
         & \quad - \frac{C_{b,i,j+\tfrac12,k}}{\Delta x} \Big(H_{z,i+\tfrac12,j+\tfrac12,k}^{n}-H_{z,i-\tfrac12,j+\tfrac12,k}^{n}\Big).
    \end{aligned}
\end{equation}

Equivalently, the tridiagonal system can be written as
\begin{align}
    \begin{split}
        -\alpha_{y1}E_{y,i,j+\tfrac12,k-1}^{n+\frac12} + \beta_{y1}E_{y,i,j+\tfrac12,k}^{n+\frac12} - \gamma_{y1}E_{y,i,j+\tfrac12,k+1}^{n+\frac12}
         & = p_{y1}E_{y,i,j+\tfrac12,k}^{n}                                                                       \\
         & \quad +\frac{1}{\Delta z}\Big(H_{x,i,j+\tfrac12,k+\tfrac12}^{n}-H_{x,i,j+\tfrac12,k-\tfrac12}^{n}\Big) \\
         & \quad -\frac{1}{\Delta x}\Big(H_{z,i+\tfrac12,j+\tfrac12,k}^{n}-H_{z,i-\tfrac12,j+\tfrac12,k}^{n}\Big) \\
         & \quad +\frac{q_{y1}}{\Delta y\Delta z}\Big(E_{z,i,j+1,k-\tfrac12}^{n}-E_{z,i,j,k-\tfrac12}^{n}\Big)    \\
         & \quad -\frac{r_{y1}}{\Delta y\Delta z}\Big(E_{z,i,j+1,k+\tfrac12}^{n}-E_{z,i,j,k+\tfrac12}^{n}\Big),
    \end{split}
    \label{eq:first-e-ey-adi-tridiagonal}
\end{align}

and the magnetic update can be written as
\begin{equation}
    H_{x,i,j+\tfrac12,k+\tfrac12}^{n+\frac12} = H_{x,i,j+\tfrac12,k+\tfrac12}^{n} + \frac{r_{y1}}{\Delta z}\Big(E_{y,i,j+\tfrac12,k+1}^{n+\frac12}-E_{y,i,j+\tfrac12,k}^{n+\frac12}\Big) - \frac{r_{y1}}{\Delta y}\Big(E_{z,i,j+1,k+\tfrac12}^{n}-E_{z,i,j,k+\tfrac12}^{n}\Big).
    \label{eq:adi-first-hx}
\end{equation}
where
\begin{align}
    \alpha_{y1} & =\frac{D_{b,i,j+\tfrac12,k-\tfrac12}}{\Delta z^2},       \\
    \beta_{y1}  & =\frac{1}{C_{b,i,j+\tfrac12,k}}+\alpha_{y1}+\gamma_{y1}, \\
    \gamma_{y1} & =\frac{D_{b,i,j+\tfrac12,k+\tfrac12}}{\Delta z^2},       \\
    p_{y1}      & =\frac{C_{a,i,j+\tfrac12,k}}{C_{b,i,j+\tfrac12,k}},      \\
    q_{y1}      & =D_{b,i,j+\tfrac12,k-\tfrac12},                          \\
    r_{y1}      & =D_{b,i,j+\tfrac12,k+\tfrac12}.
\end{align}

\paragraph{$E_z$ and $H_y$}
By equation \eqref{eq:e-z-component},
\begin{equation}
    \varepsilon_{i,j,k+\tfrac12} \frac{E_{z,i,j,k+\tfrac12}^{n+\frac12}-E_{z,i,j,k+\tfrac12}^{n}}{\Delta t/2} + \sigma_{i,j,k+\tfrac12} \frac{E_{z,i,j,k+\tfrac12}^{n+\frac12}+E_{z,i,j,k+\tfrac12}^{n}}{2} = [H_3],
    \label{eq:first-e-ez}
\end{equation}
where
\begin{equation}
    [H_3] = \frac{H_{y,i+\tfrac12,j,k+\tfrac12}^{n+\frac12}-H_{y,i-\tfrac12,j,k+\tfrac12}^{n+\frac12}}{\Delta x} - \frac{H_{x,i,j+\tfrac12,k+\tfrac12}^{n}-H_{x,i,j-\tfrac12,k+\tfrac12}^{n}}{\Delta y}.
\end{equation}
Rearranging gives
\begin{equation}
    \left(\frac{2\varepsilon}{\Delta t}+\frac{\sigma}{2}\right) E_{z}^{n+\frac12} = \left(\frac{2\varepsilon}{\Delta t}-\frac{\sigma}{2}\right) E_{z}^{n} + [H_3].
\end{equation}
Hence
\begin{equation}
    E_{z,i,j,k+\tfrac12}^{n+\frac12} = C_{a,i,j,k+\tfrac12}\,E_{z,i,j,k+\tfrac12}^{n} + C_{b,i,j,k+\tfrac12}\,[H_3],
    \label{eq:first-e-ez-adi}
\end{equation}

By equation \eqref{eq:h-y-component},
\begin{align}
    \begin{split}
        H_{y,i+\tfrac12,j,k+\tfrac12}^{n+\frac12} = H_{y,i+\tfrac12,j,k+\tfrac12}^{n} + D_{b,i+\tfrac12,j,k+\tfrac12}\Big( & \frac{E_{z,i+1,j,k+\tfrac12}^{n+\frac12}-E_{z,i,j,k+\tfrac12}^{n+\frac12}}{\Delta x} \\
                                                                                                                           & - \frac{E_{x,i+\tfrac12,j,k+1}^{n}-E_{x,i+\tfrac12,j,k}^{n}}{\Delta z}\Big),
    \end{split}
\end{align}

Substituting the magnetic update into \eqref{eq:first-e-ez-adi} and eliminating $H_{y}^{n+\frac12}$ terms gives
\begin{align}
    \begin{split}
        E_{z,i,j,k+\tfrac12}^{n+\frac12} & = C_{a,i,j,k+\tfrac12}\,E_{z,i,j,k+\tfrac12}^{n}                                                                                                                                                                                                                                                           \\
                                         & \quad + \frac{C_{b,i,j,k+\tfrac12}}{\Delta x}\Bigg[H_{y,i+\tfrac12,j,k+\tfrac12}^{n} + D_{b,i+\tfrac12,j,k+\tfrac12}\left(\frac{E_{z,i+1,j,k+\tfrac12}^{n+\frac12}-E_{z,i,j,k+\tfrac12}^{n+\frac12}}{\Delta x}\right. \left. - \frac{E_{x,i+\tfrac12,j,k+1}^{n}-E_{x,i+\tfrac12,j,k}^{n}}{\Delta z}\right) \\
                                         & \qquad\qquad\qquad - H_{y,i-\tfrac12,j,k+\tfrac12}^{n} - D_{b,i-\tfrac12,j,k+\tfrac12}\left(\frac{E_{z,i,j,k+\tfrac12}^{n+\frac12}-E_{z,i-1,j,k+\tfrac12}^{n+\frac12}}{\Delta x}\right. \left. - \frac{E_{x,i-\tfrac12,j,k+1}^{n}-E_{x,i-\tfrac12,j,k}^{n}}{\Delta z}\right)\Bigg]                         \\
                                         & \quad - \frac{C_{b,i,j,k+\tfrac12}}{\Delta y}\left(H_{x,i,j+\tfrac12,k+\tfrac12}^{n}-H_{x,i,j-\tfrac12,k+\tfrac12}^{n}\right).
    \end{split}
\end{align}

Collecting the unknown $E_{z}^{n+\frac12}$ terms gives
\begin{equation}
    \begin{aligned}
         & \quad -\frac{C_{b,i,j,k+\tfrac12}\,D_{b,i+\tfrac12,j,k+\tfrac12}}{\Delta x^2}\, E_{z,i+1,j,k+\tfrac12}^{n+\frac12}                                                                                     \\
         & \quad + \left[1+\frac{C_{b,i,j,k+\tfrac12}\,D_{b,i+\tfrac12,j,k+\tfrac12}}{\Delta x^2} +\frac{C_{b,i,j,k+\tfrac12}\,D_{b,i-\tfrac12,j,k+\tfrac12}}{\Delta x^2}\right] E_{z,i,j,k+\tfrac12}^{n+\frac12} \\
         & \quad - \frac{C_{b,i,j,k+\tfrac12}\,D_{b,i-\tfrac12,j,k+\tfrac12}}{\Delta x^2}\, E_{z,i-1,j,k+\tfrac12}^{n+\frac12}                                                                                    \\
         & = C_{a,i,j,k+\tfrac12}\, E_{z,i,j,k+\tfrac12}^{n} + \frac{C_{b,i,j,k+\tfrac12}}{\Delta x} \Big(H_{y,i+\tfrac12,j,k+\tfrac12}^{n}-H_{y,i-\tfrac12,j,k+\tfrac12}^{n}\Big)                                \\
         & \quad - \frac{C_{b,i,j,k+\tfrac12}\,D_{b,i+\tfrac12,j,k+\tfrac12}}{\Delta x\,\Delta z} \Big(E_{x,i+\tfrac12,j,k+1}^{n}-E_{x,i+\tfrac12,j,k}^{n}\Big)                                                   \\
         & \quad + \frac{C_{b,i,j,k+\tfrac12}\,D_{b,i-\tfrac12,j,k+\tfrac12}}{\Delta x\,\Delta z} \Big(E_{x,i-\tfrac12,j,k+1}^{n}-E_{x,i-\tfrac12,j,k}^{n}\Big)                                                   \\
         & \quad - \frac{C_{b,i,j,k+\tfrac12}}{\Delta y} \Big(H_{x,i,j+\tfrac12,k+\tfrac12}^{n}-H_{x,i,j-\tfrac12,k+\tfrac12}^{n}\Big).
    \end{aligned}
\end{equation}

Equivalently, the tridiagonal system can be written as
\begin{align}
    \begin{split}
        -\alpha_{z1}E_{z,i-1,j,k+\tfrac12}^{n+\frac12} + \beta_{z1}E_{z,i,j,k+\tfrac12}^{n+\frac12} - \gamma_{z1}E_{z,i+1,j,k+\tfrac12}^{n+\frac12}
         & = p_{z1}E_{z,i,j,k+\tfrac12}^{n}                                                                       \\
         & \quad +\frac{1}{\Delta x}\Big(H_{y,i+\tfrac12,j,k+\tfrac12}^{n}-H_{y,i-\tfrac12,j,k+\tfrac12}^{n}\Big) \\
         & \quad -\frac{1}{\Delta y}\Big(H_{x,i,j+\tfrac12,k+\tfrac12}^{n}-H_{x,i,j-\tfrac12,k+\tfrac12}^{n}\Big) \\
         & \quad +\frac{q_{z1}}{\Delta z\Delta x}\Big(E_{x,i-\tfrac12,j,k+1}^{n}-E_{x,i-\tfrac12,j,k}^{n}\Big)    \\
         & \quad -\frac{r_{z1}}{\Delta z\Delta x}\Big(E_{x,i+\tfrac12,j,k+1}^{n}-E_{x,i+\tfrac12,j,k}^{n}\Big),
    \end{split}
    \label{eq:first-e-ez-adi-tridiagonal}
\end{align}

and the magnetic update can be written as
\begin{equation}
    H_{y,i+\tfrac12,j,k+\tfrac12}^{n+\frac12} = H_{y,i+\tfrac12,j,k+\tfrac12}^{n} + \frac{r_{z1}}{\Delta x}\Big(E_{z,i+1,j,k+\tfrac12}^{n+\frac12}-E_{z,i,j,k+\tfrac12}^{n+\frac12}\Big) - \frac{r_{z1}}{\Delta z}\Big(E_{x,i+\tfrac12,j,k+1}^{n}-E_{x,i+\tfrac12,j,k}^{n}\Big).
    \label{eq:adi-first-hy}
\end{equation}
where
\begin{align}
    \alpha_{z1} & =\frac{D_{b,i-\tfrac12,j,k+\tfrac12}}{\Delta x^2},       \\
    \beta_{z1}  & =\frac{1}{C_{b,i,j,k+\tfrac12}}+\alpha_{z1}+\gamma_{z1}, \\
    \gamma_{z1} & =\frac{D_{b,i+\tfrac12,j,k+\tfrac12}}{\Delta x^2},       \\
    p_{z1}      & =\frac{C_{a,i,j,k+\tfrac12}}{C_{b,i,j,k+\tfrac12}},      \\
    q_{z1}      & =D_{b,i-\tfrac12,j,k+\tfrac12},                          \\
    r_{z1}      & =D_{b,i+\tfrac12,j,k+\tfrac12}.
\end{align}

In the AMReX Yee grid, with vacuum media, the field updates are implemented as
\begin{subequations}\label{eq:adi-first-half-amrex}
    \begin{align}
        \begin{split}
            - E_x^{n+\frac12}(i, j - 1, k) + \left(2+\frac{4\Delta y^2}{c^2\Delta t^2}\right) E_x^{n+\frac12}(i, j, k) & - E_x^{n+\frac12}(i, j + 1, k) = \frac{4\Delta y^2}{c^2\Delta t^2} E_x^{n}(i, j, k)                                                                        \\
                                                                                                                       & + \frac{2\Delta y^2}{\Delta t}\left[\frac{B_z^{n}(i, j, k)-B_z^{n}(i, j - 1, k)}{\Delta y} - \frac{B_y^{n}(i, j, k)-B_y^{n}(i, j, k - 1)}{\Delta z}\right] \\
                                                                                                                       & + \Delta y\left[\frac{E_y^{n}(i + 1, j - 1, k)-E_y^{n}(i, j - 1, k)}{\Delta x} - \frac{E_y^{n}(i + 1, j, k)-E_y^{n}(i, j, k)}{\Delta x}\right],            \\
        \end{split}
    \end{align}
    \begin{align}
        \begin{split}
            - E_y^{n+\frac12}(i, j, k - 1) + \left(2+\frac{4\Delta z^2}{c^2\Delta t^2}\right) E_y^{n+\frac12}(i, j, k) & - E_y^{n+\frac12}(i, j, k + 1) = \frac{4\Delta z^2}{c^2\Delta t^2} E_y^{n}(i, j, k)                                                                        \\
                                                                                                                       & + \frac{2\Delta z^2}{\Delta t}\left[\frac{B_x^{n}(i, j, k)-B_x^{n}(i, j, k - 1)}{\Delta z} - \frac{B_z^{n}(i, j, k)-B_z^{n}(i - 1, j, k)}{\Delta x}\right] \\
                                                                                                                       & + \Delta z\left[\frac{E_z^{n}(i, j + 1, k - 1)-E_z^{n}(i, j, k - 1)}{\Delta y} - \frac{E_z^{n}(i, j + 1, k)-E_z^{n}(i, j, k)}{\Delta y}\right],            \\
        \end{split}
    \end{align}
    \begin{align}
        \begin{split}
            - E_z^{n+\frac12}(i - 1, j, k) + \left(2+\frac{4\Delta x^2}{c^2\Delta t^2}\right) E_z^{n+\frac12}(i, j, k) & - E_z^{n+\frac12}(i + 1, j, k) = \frac{4\Delta x^2}{c^2\Delta t^2} E_z^{n}(i, j, k)                                                                        \\
                                                                                                                       & + \frac{2\Delta x^2}{\Delta t}\left[\frac{B_y^{n}(i, j, k)-B_y^{n}(i - 1, j, k)}{\Delta x} - \frac{B_x^{n}(i, j, k)-B_x^{n}(i, j - 1, k)}{\Delta y}\right] \\
                                                                                                                       & + \Delta x\left[\frac{E_x^{n}(i - 1, j, k + 1)-E_x^{n}(i - 1, j, k)}{\Delta z} - \frac{E_x^{n}(i, j, k + 1)-E_x^{n}(i, j, k)}{\Delta z}\right],            \\
        \end{split}
    \end{align}
    \begin{align}
        B_x^{n+\frac12}(i, j, k)
         & = B_x^{n}(i, j, k) + \frac{\Delta t}{2}\left[\frac{E_y^{n+\frac12}(i, j, k + 1)-E_y^{n+\frac12}(i, j, k)}{\Delta z} - \frac{E_z^{n}(i, j + 1, k)-E_z^{n}(i, j, k)}{\Delta y}\right], \\
        B_y^{n+\frac12}(i, j, k)
         & = B_y^{n}(i, j, k) + \frac{\Delta t}{2}\left[\frac{E_z^{n+\frac12}(i + 1, j, k)-E_z^{n+\frac12}(i, j, k)}{\Delta x} - \frac{E_x^{n}(i, j, k + 1)-E_x^{n}(i, j, k)}{\Delta z}\right], \\
        B_z^{n+\frac12}(i, j, k)
         & = B_z^{n}(i, j, k) + \frac{\Delta t}{2}\left[\frac{E_x^{n+\frac12}(i, j + 1, k)-E_x^{n+\frac12}(i, j, k)}{\Delta y} - \frac{E_y^{n}(i + 1, j, k)-E_y^{n}(i, j, k)}{\Delta x}\right].
    \end{align}
\end{subequations}

\subsection{Second ADI Half-Step}

\paragraph{$E_x$ and $H_y$}
By equation \eqref{eq:e-x-component},
\begin{equation}
    \varepsilon_{i+\tfrac12,j,k} \frac{E_{x,i+\tfrac12,j,k}^{n+1}-E_{x,i+\tfrac12,j,k}^{n+\frac12}}{\Delta t/2} + \sigma_{i+\tfrac12,j,k} \frac{E_{x,i+\tfrac12,j,k}^{n+1}+E_{x,i+\tfrac12,j,k}^{n+\frac12}}{2} = [H_4],
    \label{eq:second-e-ex}
\end{equation}
where
\begin{equation}
    [H_4] = \frac{H_{z,i+\tfrac12,j+\tfrac12,k}^{n+\frac12}-H_{z,i+\tfrac12,j-\tfrac12,k}^{n+\frac12}}{\Delta y} - \frac{H_{y,i+\tfrac12,j,k+\tfrac12}^{n+1}-H_{y,i+\tfrac12,j,k-\tfrac12}^{n+1}}{\Delta z}.
\end{equation}
Rearranging gives
\begin{equation}
    \left(\frac{2\varepsilon}{\Delta t}+\frac{\sigma}{2}\right) E_{x}^{n+1} = \left(\frac{2\varepsilon}{\Delta t}-\frac{\sigma}{2}\right) E_{x}^{n+\frac12} + [H_4].
\end{equation}
Hence
\begin{equation}
    E_{x,i+\tfrac12,j,k}^{n+1} = C_{a,i+\tfrac12,j,k}\,E_{x,i+\tfrac12,j,k}^{n+\frac12} + C_{b,i+\tfrac12,j,k}\,[H_4],
    \label{eq:second-e-ex-adi}
\end{equation}

By equation \eqref{eq:h-y-component},
\begin{align}
    \begin{split}
        H_{y,i+\tfrac12,j,k+\tfrac12}^{n+1} = H_{y,i+\tfrac12,j,k+\tfrac12}^{n+\frac12} + D_{b,i+\tfrac12,j,k+\tfrac12}\Big( & \frac{E_{z,i+1,j,k+\tfrac12}^{n+\frac12}-E_{z,i,j,k+\tfrac12}^{n+\frac12}}{\Delta x} \\
                                                                                                                             & - \frac{E_{x,i+\tfrac12,j,k+1}^{n+1}-E_{x,i+\tfrac12,j,k}^{n+1}}{\Delta z}\Big),
    \end{split}
\end{align}

Substituting the magnetic update into \eqref{eq:second-e-ex-adi} and eliminating $H_{y}^{n+1}$ terms gives
\begin{align}
    \begin{split}
        E_{x,i+\tfrac12,j,k}^{n+1} & = C_{a,i+\tfrac12,j,k}\,E_{x,i+\tfrac12,j,k}^{n+\frac12}                                                                                                                                                                                                                                                               \\
                                   & \quad + \frac{C_{b,i+\tfrac12,j,k}}{\Delta y}\Big(H_{z,i+\tfrac12,j+\tfrac12,k}^{n+\frac12}-H_{z,i+\tfrac12,j-\tfrac12,k}^{n+\frac12}\Big)                                                                                                                                                                             \\
                                   & \quad - \frac{C_{b,i+\tfrac12,j,k}}{\Delta z}\Bigg[H_{y,i+\tfrac12,j,k+\tfrac12}^{n+\frac12} + D_{b,i+\tfrac12,j,k+\tfrac12}\left(\frac{E_{z,i+1,j,k+\tfrac12}^{n+\frac12}-E_{z,i,j,k+\tfrac12}^{n+\frac12}}{\Delta x}\right. \left. - \frac{E_{x,i+\tfrac12,j,k+1}^{n+1}-E_{x,i+\tfrac12,j,k}^{n+1}}{\Delta z}\right) \\
                                   & \qquad\qquad\qquad - H_{y,i+\tfrac12,j,k-\tfrac12}^{n+\frac12} - D_{b,i+\tfrac12,j,k-\tfrac12}\left(\frac{E_{z,i+1,j,k-\tfrac12}^{n+\frac12}-E_{z,i,j,k-\tfrac12}^{n+\frac12}}{\Delta x}\right. \left. - \frac{E_{x,i+\tfrac12,j,k}^{n+1}-E_{x,i+\tfrac12,j,k-1}^{n+1}}{\Delta z}\right)\Bigg].
    \end{split}
\end{align}

Collecting the unknown $E_{x}^{n+1}$ terms gives
\begin{equation}
    \begin{aligned}
         & \quad -\frac{C_{b,i+\tfrac12,j,k}\,D_{b,i+\tfrac12,j,k+\tfrac12}}{\Delta z^2}\, E_{x,i+\tfrac12,j,k+1}^{n+1}                                                                                     \\
         & \quad + \left[1+\frac{C_{b,i+\tfrac12,j,k}\,D_{b,i+\tfrac12,j,k+\tfrac12}}{\Delta z^2} +\frac{C_{b,i+\tfrac12,j,k}\,D_{b,i+\tfrac12,j,k-\tfrac12}}{\Delta z^2}\right] E_{x,i+\tfrac12,j,k}^{n+1} \\
         & \quad - \frac{C_{b,i+\tfrac12,j,k}\,D_{b,i+\tfrac12,j,k-\tfrac12}}{\Delta z^2}\, E_{x,i+\tfrac12,j,k-1}^{n+1}                                                                                    \\
         & = C_{a,i+\tfrac12,j,k}\, E_{x,i+\tfrac12,j,k}^{n+\frac12} + \frac{C_{b,i+\tfrac12,j,k}}{\Delta y} \Big(H_{z,i+\tfrac12,j+\tfrac12,k}^{n+\frac12}-H_{z,i+\tfrac12,j-\tfrac12,k}^{n+\frac12}\Big)  \\
         & \quad - \frac{C_{b,i+\tfrac12,j,k}}{\Delta z} \Big(H_{y,i+\tfrac12,j,k+\tfrac12}^{n+\frac12}-H_{y,i+\tfrac12,j,k-\tfrac12}^{n+\frac12}\Big)                                                      \\
         & \quad - \frac{C_{b,i+\tfrac12,j,k}\,D_{b,i+\tfrac12,j,k+\tfrac12}}{\Delta z\,\Delta x} \Big(E_{z,i+1,j,k+\tfrac12}^{n+\frac12}-E_{z,i,j,k+\tfrac12}^{n+\frac12}\Big)                             \\
         & \quad + \frac{C_{b,i+\tfrac12,j,k}\,D_{b,i+\tfrac12,j,k-\tfrac12}}{\Delta z\,\Delta x} \Big(E_{z,i+1,j,k-\tfrac12}^{n+\frac12}-E_{z,i,j,k-\tfrac12}^{n+\frac12}\Big).
    \end{aligned}
\end{equation}

Equivalently, the tridiagonal system can be written as
\begin{align}
    \begin{split}
        -\alpha_{x2}E_{x,i+\tfrac12,j,k-1}^{n+1} + \beta_{x2}E_{x,i+\tfrac12,j,k}^{n+1} - \gamma_{x2}E_{x,i+\tfrac12,j,k+1}^{n+1}
         & = p_{x2}E_{x,i+\tfrac12,j,k}^{n+\frac12}                                                                               \\
         & \quad +\frac{1}{\Delta y}\Big(H_{z,i+\tfrac12,j+\tfrac12,k}^{n+\frac12}-H_{z,i+\tfrac12,j-\tfrac12,k}^{n+\frac12}\Big) \\
         & \quad -\frac{1}{\Delta z}\Big(H_{y,i+\tfrac12,j,k+\tfrac12}^{n+\frac12}-H_{y,i+\tfrac12,j,k-\tfrac12}^{n+\frac12}\Big) \\
         & \quad +\frac{q_{x2}}{\Delta x\Delta z}\Big(E_{z,i+1,j,k-\tfrac12}^{n+\frac12}-E_{z,i,j,k-\tfrac12}^{n+\frac12}\Big)    \\
         & \quad -\frac{r_{x2}}{\Delta x\Delta z}\Big(E_{z,i+1,j,k+\tfrac12}^{n+\frac12}-E_{z,i,j,k+\tfrac12}^{n+\frac12}\Big),
    \end{split}
    \label{eq:second-e-ex-adi-tridiagonal}
\end{align}

and the magnetic update can be written as
\begin{equation}
    H_{y,i+\tfrac12,j,k+\tfrac12}^{n+1} = H_{y,i+\tfrac12,j,k+\tfrac12}^{n+\frac12} + \frac{r_{x2}}{\Delta x}\Big(E_{z,i+1,j,k+\tfrac12}^{n+\frac12}-E_{z,i,j,k+\tfrac12}^{n+\frac12}\Big) - \frac{r_{x2}}{\Delta z}\Big(E_{x,i+\tfrac12,j,k+1}^{n+1}-E_{x,i+\tfrac12,j,k}^{n+1}\Big).
    \label{eq:adi-second-hy}
\end{equation}
where
\begin{align}
    \alpha_{x2} & =\frac{D_{b,i+\tfrac12,j,k-\tfrac12}}{\Delta z^2},       \\
    \beta_{x2}  & =\frac{1}{C_{b,i+\tfrac12,j,k}}+\alpha_{x2}+\gamma_{x2}, \\
    \gamma_{x2} & =\frac{D_{b,i+\tfrac12,j,k+\tfrac12}}{\Delta z^2},       \\
    p_{x2}      & =\frac{C_{a,i+\tfrac12,j,k}}{C_{b,i+\tfrac12,j,k}},      \\
    q_{x2}      & =D_{b,i+\tfrac12,j,k-\tfrac12},                          \\
    r_{x2}      & =D_{b,i+\tfrac12,j,k+\tfrac12}.
\end{align}



\paragraph{$E_y$ and $H_z$}
By equation \eqref{eq:e-y-component},
\begin{equation}
    \varepsilon_{i,j+\tfrac12,k} \frac{E_{y,i,j+\tfrac12,k}^{n+1}-E_{y,i,j+\tfrac12,k}^{n+\frac12}}{\Delta t/2} + \sigma_{i,j+\tfrac12,k} \frac{E_{y,i,j+\tfrac12,k}^{n+1}+E_{y,i,j+\tfrac12,k}^{n+\frac12}}{2} = [H_5],
    \label{eq:second-e-ey}
\end{equation}
where
\begin{equation}
    [H_5] = \frac{H_{x,i,j+\tfrac12,k+\tfrac12}^{n+\frac12}-H_{x,i,j+\tfrac12,k-\tfrac12}^{n+\frac12}}{\Delta z} - \frac{H_{z,i+\tfrac12,j+\tfrac12,k}^{n+1}-H_{z,i-\tfrac12,j+\tfrac12,k}^{n+1}}{\Delta x}.
\end{equation}
Rearranging gives
\begin{equation}
    \left(\frac{2\varepsilon}{\Delta t}+\frac{\sigma}{2}\right) E_{y}^{n+1} = \left(\frac{2\varepsilon}{\Delta t}-\frac{\sigma}{2}\right) E_{y}^{n+\frac12} + [H_5].
\end{equation}
Hence
\begin{equation}
    E_{y,i,j+\tfrac12,k}^{n+1} = C_{a,i,j+\tfrac12,k}\,E_{y,i,j+\tfrac12,k}^{n+\frac12} + C_{b,i,j+\tfrac12,k}\,[H_5],
    \label{eq:second-e-ey-adi}
\end{equation}

By equation \eqref{eq:h-z-component},
\begin{align}
    \begin{split}
        H_{z,i+\tfrac12,j+\tfrac12,k}^{n+1} = H_{z,i+\tfrac12,j+\tfrac12,k}^{n+\frac12} + D_{b,i+\tfrac12,j+\tfrac12,k}\Big( & \frac{E_{x,i+\tfrac12,j+1,k}^{n+\frac12}-E_{x,i+\tfrac12,j,k}^{n+\frac12}}{\Delta y} \\
                                                                                                                             & - \frac{E_{y,i+1,j+\tfrac12,k}^{n+1}-E_{y,i,j+\tfrac12,k}^{n+1}}{\Delta x}\Big),
    \end{split}
\end{align}

Substituting the magnetic update into \eqref{eq:second-e-ey-adi} and eliminating $H_{z}^{n+1}$ terms gives
\begin{align}
    \begin{split}
        E_{y,i,j+\tfrac12,k}^{n+1} & = C_{a,i,j+\tfrac12,k}\,E_{y,i,j+\tfrac12,k}^{n+\frac12}                                                                                                                                                                                                                                                               \\
                                   & \quad + \frac{C_{b,i,j+\tfrac12,k}}{\Delta z}\Big(H_{x,i,j+\tfrac12,k+\tfrac12}^{n+\frac12}-H_{x,i,j+\tfrac12,k-\tfrac12}^{n+\frac12}\Big)                                                                                                                                                                             \\
                                   & \quad - \frac{C_{b,i,j+\tfrac12,k}}{\Delta x}\Bigg[H_{z,i+\tfrac12,j+\tfrac12,k}^{n+\frac12} + D_{b,i+\tfrac12,j+\tfrac12,k}\left(\frac{E_{x,i+\tfrac12,j+1,k}^{n+\frac12}-E_{x,i+\tfrac12,j,k}^{n+\frac12}}{\Delta y}\right. \left. - \frac{E_{y,i+1,j+\tfrac12,k}^{n+1}-E_{y,i,j+\tfrac12,k}^{n+1}}{\Delta x}\right) \\
                                   & \qquad\qquad\qquad - H_{z,i-\tfrac12,j+\tfrac12,k}^{n+\frac12} - D_{b,i-\tfrac12,j+\tfrac12,k}\left(\frac{E_{x,i-\tfrac12,j+1,k}^{n+\frac12}-E_{x,i-\tfrac12,j,k}^{n+\frac12}}{\Delta y}\right. \left. - \frac{E_{y,i,j+\tfrac12,k}^{n+1}-E_{y,i-1,j+\tfrac12,k}^{n+1}}{\Delta x}\right)\Bigg].
    \end{split}
\end{align}

Collecting the unknown $E_{y}^{n+1}$ terms gives
\begin{equation}
    \begin{aligned}
         & \quad -\frac{C_{b,i,j+\tfrac12,k}\,D_{b,i+\tfrac12,j+\tfrac12,k}}{\Delta x^2}\, E_{y,i+1,j+\tfrac12,k}^{n+1}                                                                                     \\
         & \quad + \left[1+\frac{C_{b,i,j+\tfrac12,k}\,D_{b,i+\tfrac12,j+\tfrac12,k}}{\Delta x^2} +\frac{C_{b,i,j+\tfrac12,k}\,D_{b,i-\tfrac12,j+\tfrac12,k}}{\Delta x^2}\right] E_{y,i,j+\tfrac12,k}^{n+1} \\
         & \quad - \frac{C_{b,i,j+\tfrac12,k}\,D_{b,i-\tfrac12,j+\tfrac12,k}}{\Delta x^2}\, E_{y,i-1,j+\tfrac12,k}^{n+1}                                                                                    \\
         & = C_{a,i,j+\tfrac12,k}\, E_{y,i,j+\tfrac12,k}^{n+\frac12} + \frac{C_{b,i,j+\tfrac12,k}}{\Delta z} \Big(H_{x,i,j+\tfrac12,k+\tfrac12}^{n+\frac12}-H_{x,i,j+\tfrac12,k-\tfrac12}^{n+\frac12}\Big)  \\
         & \quad - \frac{C_{b,i,j+\tfrac12,k}}{\Delta x} \Big(H_{z,i+\tfrac12,j+\tfrac12,k}^{n+\frac12}-H_{z,i-\tfrac12,j+\tfrac12,k}^{n+\frac12}\Big)                                                      \\
         & \quad - \frac{C_{b,i,j+\tfrac12,k}\,D_{b,i+\tfrac12,j+\tfrac12,k}}{\Delta x\,\Delta y} \Big(E_{x,i+\tfrac12,j+1,k}^{n+\frac12}-E_{x,i+\tfrac12,j,k}^{n+\frac12}\Big)                             \\
         & \quad + \frac{C_{b,i,j+\tfrac12,k}\,D_{b,i-\tfrac12,j+\tfrac12,k}}{\Delta x\,\Delta y} \Big(E_{x,i-\tfrac12,j+1,k}^{n+\frac12}-E_{x,i-\tfrac12,j,k}^{n+\frac12}\Big).
    \end{aligned}
\end{equation}

Equivalently, the tridiagonal system can be written as
\begin{align}
    \begin{split}
        -\alpha_{y2}E_{y,i-1,j+\tfrac12,k}^{n+1} + \beta_{y2}E_{y,i,j+\tfrac12,k}^{n+1} - \gamma_{y2}E_{y,i+1,j+\tfrac12,k}^{n+1}
         & = p_{y2}E_{y,i,j+\tfrac12,k}^{n+\frac12}                                                                               \\
         & \quad +\frac{1}{\Delta z}\Big(H_{x,i,j+\tfrac12,k+\tfrac12}^{n+\frac12}-H_{x,i,j+\tfrac12,k-\tfrac12}^{n+\frac12}\Big) \\
         & \quad -\frac{1}{\Delta x}\Big(H_{z,i+\tfrac12,j+\tfrac12,k}^{n+\frac12}-H_{z,i-\tfrac12,j+\tfrac12,k}^{n+\frac12}\Big) \\
         & \quad +\frac{q_{y2}}{\Delta y\Delta x}\Big(E_{x,i-\tfrac12,j+1,k}^{n+\frac12}-E_{x,i-\tfrac12,j,k}^{n+\frac12}\Big)    \\
         & \quad -\frac{r_{y2}}{\Delta y\Delta x}\Big(E_{x,i+\tfrac12,j+1,k}^{n+\frac12}-E_{x,i+\tfrac12,j,k}^{n+\frac12}\Big),
    \end{split}
    \label{eq:second-e-ey-adi-tridiagonal}
\end{align}

and the magnetic update can be written as
\begin{equation}
    H_{z,i+\tfrac12,j+\tfrac12,k}^{n+1} = H_{z,i+\tfrac12,j+\tfrac12,k}^{n+\frac12} + \frac{r_{y2}}{\Delta y}\Big(E_{x,i+\tfrac12,j+1,k}^{n+\frac12}-E_{x,i+\tfrac12,j,k}^{n+\frac12}\Big) - \frac{r_{y2}}{\Delta x}\Big(E_{y,i+1,j+\tfrac12,k}^{n+1}-E_{y,i,j+\tfrac12,k}^{n+1}\Big).
    \label{eq:adi-second-hz}
\end{equation}
where
\begin{align}
    \alpha_{y2} & =\frac{D_{b,i-\tfrac12,j+\tfrac12,k}}{\Delta x^2},       \\
    \beta_{y2}  & =\frac{1}{C_{b,i,j+\tfrac12,k}}+\alpha_{y2}+\gamma_{y2}, \\
    \gamma_{y2} & =\frac{D_{b,i+\tfrac12,j+\tfrac12,k}}{\Delta x^2},       \\
    p_{y2}      & =\frac{C_{a,i,j+\tfrac12,k}}{C_{b,i,j+\tfrac12,k}},      \\
    q_{y2}      & =D_{b,i-\tfrac12,j+\tfrac12,k},                          \\
    r_{y2}      & =D_{b,i+\tfrac12,j+\tfrac12,k}.
\end{align}

\paragraph{$E_z$ and $H_x$}
By equation \eqref{eq:e-z-component},
\begin{equation}
    \varepsilon_{i,j,k+\tfrac12} \frac{E_{z,i,j,k+\tfrac12}^{n+1}-E_{z,i,j,k+\tfrac12}^{n+\frac12}}{\Delta t/2} + \sigma_{i,j,k+\tfrac12} \frac{E_{z,i,j,k+\tfrac12}^{n+1}+E_{z,i,j,k+\tfrac12}^{n+\frac12}}{2} = [H_6],
    \label{eq:second-e-ez}
\end{equation}
where
\begin{equation}
    [H_6] = \frac{H_{y,i+\tfrac12,j,k+\tfrac12}^{n+\frac12}-H_{y,i-\tfrac12,j,k+\tfrac12}^{n+\frac12}}{\Delta x} - \frac{H_{x,i,j+\tfrac12,k+\tfrac12}^{n+1}-H_{x,i,j-\tfrac12,k+\tfrac12}^{n+1}}{\Delta y}.
\end{equation}
Rearranging gives
\begin{equation}
    \left(\frac{2\varepsilon}{\Delta t}+\frac{\sigma}{2}\right) E_{z}^{n+1} = \left(\frac{2\varepsilon}{\Delta t}-\frac{\sigma}{2}\right) E_{z}^{n+\frac12} + [H_6].
\end{equation}
Hence
\begin{equation}
    E_{z,i,j,k+\tfrac12}^{n+1} = C_{a,i,j,k+\tfrac12}\,E_{z,i,j,k+\tfrac12}^{n+\frac12} + C_{b,i,j,k+\tfrac12}\,[H_6],
    \label{eq:second-e-ez-adi}
\end{equation}

By equation \eqref{eq:h-x-component},
\begin{align}
    \begin{split}
        H_{x,i,j+\tfrac12,k+\tfrac12}^{n+1} = H_{x,i,j+\tfrac12,k+\tfrac12}^{n+\frac12} + D_{b,i,j+\tfrac12,k+\tfrac12}\Big( & \frac{E_{y,i,j+\tfrac12,k+1}^{n+\frac12}-E_{y,i,j+\tfrac12,k}^{n+\frac12}}{\Delta z} \\
                                                                                                                             & - \frac{E_{z,i,j+1,k+\tfrac12}^{n+1}-E_{z,i,j,k+\tfrac12}^{n+1}}{\Delta y}\Big),
    \end{split}
\end{align}

Substituting the magnetic update into \eqref{eq:second-e-ez-adi} and eliminating $H_{x}^{n+1}$ terms gives
\begin{align}
    \begin{split}
        E_{z,i,j,k+\tfrac12}^{n+1} & = C_{a,i,j,k+\tfrac12}\,E_{z,i,j,k+\tfrac12}^{n+\frac12}                                                                                                                                                                                                                                                               \\
                                   & \quad + \frac{C_{b,i,j,k+\tfrac12}}{\Delta x}\Big(H_{y,i+\tfrac12,j,k+\tfrac12}^{n+\frac12}-H_{y,i-\tfrac12,j,k+\tfrac12}^{n+\frac12}\Big)                                                                                                                                                                             \\
                                   & \quad - \frac{C_{b,i,j,k+\tfrac12}}{\Delta y}\Bigg[H_{x,i,j+\tfrac12,k+\tfrac12}^{n+\frac12} + D_{b,i,j+\tfrac12,k+\tfrac12}\left(\frac{E_{y,i,j+\tfrac12,k+1}^{n+\frac12}-E_{y,i,j+\tfrac12,k}^{n+\frac12}}{\Delta z}\right. \left. - \frac{E_{z,i,j+1,k+\tfrac12}^{n+1}-E_{z,i,j,k+\tfrac12}^{n+1}}{\Delta y}\right) \\
                                   & \qquad\qquad\qquad - H_{x,i,j-\tfrac12,k+\tfrac12}^{n+\frac12} - D_{b,i,j-\tfrac12,k+\tfrac12}\left(\frac{E_{y,i,j-\tfrac12,k+1}^{n+\frac12}-E_{y,i,j-\tfrac12,k}^{n+\frac12}}{\Delta z}\right. \left. - \frac{E_{z,i,j,k+\tfrac12}^{n+1}-E_{z,i,j-1,k+\tfrac12}^{n+1}}{\Delta y}\right)\Bigg].
    \end{split}
\end{align}

Collecting the unknown $E_{z}^{n+1}$ terms gives
\begin{equation}
    \begin{aligned}
         & \quad -\frac{C_{b,i,j,k+\tfrac12}\,D_{b,i,j+\tfrac12,k+\tfrac12}}{\Delta y^2}\, E_{z,i,j+1,k+\tfrac12}^{n+1}                                                                                     \\
         & \quad + \left[1+\frac{C_{b,i,j,k+\tfrac12}\,D_{b,i,j+\tfrac12,k+\tfrac12}}{\Delta y^2} +\frac{C_{b,i,j,k+\tfrac12}\,D_{b,i,j-\tfrac12,k+\tfrac12}}{\Delta y^2}\right] E_{z,i,j,k+\tfrac12}^{n+1} \\
         & \quad - \frac{C_{b,i,j,k+\tfrac12}\,D_{b,i,j-\tfrac12,k+\tfrac12}}{\Delta y^2}\, E_{z,i,j-1,k+\tfrac12}^{n+1}                                                                                    \\
         & = C_{a,i,j,k+\tfrac12}\, E_{z,i,j,k+\tfrac12}^{n+\frac12} + \frac{C_{b,i,j,k+\tfrac12}}{\Delta x} \Big(H_{y,i+\tfrac12,j,k+\tfrac12}^{n+\frac12}-H_{y,i-\tfrac12,j,k+\tfrac12}^{n+\frac12}\Big)  \\
         & \quad - \frac{C_{b,i,j,k+\tfrac12}}{\Delta y} \Big(H_{x,i,j+\tfrac12,k+\tfrac12}^{n+\frac12}-H_{x,i,j-\tfrac12,k+\tfrac12}^{n+\frac12}\Big)                                                      \\
         & \quad - \frac{C_{b,i,j,k+\tfrac12}\,D_{b,i,j+\tfrac12,k+\tfrac12}}{\Delta y\,\Delta z} \Big(E_{y,i,j+\tfrac12,k+1}^{n+\frac12}-E_{y,i,j+\tfrac12,k}^{n+\frac12}\Big)                             \\
         & \quad + \frac{C_{b,i,j,k+\tfrac12}\,D_{b,i,j-\tfrac12,k+\tfrac12}}{\Delta y\,\Delta z} \Big(E_{y,i,j-\tfrac12,k+1}^{n+\frac12}-E_{y,i,j-\tfrac12,k}^{n+\frac12}\Big).
    \end{aligned}
\end{equation}

Equivalently, the tridiagonal system can be written as
\begin{align}
    \begin{split}
        -\alpha_{z2}E_{z,i,j-1,k+\tfrac12}^{n+1} + \beta_{z2}E_{z,i,j,k+\tfrac12}^{n+1} - \gamma_{z2}E_{z,i,j+1,k+\tfrac12}^{n+1}
         & = p_{z2}E_{z,i,j,k+\tfrac12}^{n+\frac12}                                                                               \\
         & \quad +\frac{1}{\Delta x}\Big(H_{y,i+\tfrac12,j,k+\tfrac12}^{n+\frac12}-H_{y,i-\tfrac12,j,k+\tfrac12}^{n+\frac12}\Big) \\
         & \quad -\frac{1}{\Delta y}\Big(H_{x,i,j+\tfrac12,k+\tfrac12}^{n+\frac12}-H_{x,i,j-\tfrac12,k+\tfrac12}^{n+\frac12}\Big) \\
         & \quad +\frac{q_{z2}}{\Delta z\Delta y}\Big(E_{y,i,j-\tfrac12,k+1}^{n+\frac12}-E_{y,i,j-\tfrac12,k}^{n+\frac12}\Big)    \\
         & \quad -\frac{r_{z2}}{\Delta z\Delta y}\Big(E_{y,i,j+\tfrac12,k+1}^{n+\frac12}-E_{y,i,j+\tfrac12,k}^{n+\frac12}\Big),
    \end{split}
    \label{eq:second-e-ez-adi-tridiagonal}
\end{align}

and the magnetic update can be written as
\begin{equation}
    H_{x,i,j+\tfrac12,k+\tfrac12}^{n+1} = H_{x,i,j+\tfrac12,k+\tfrac12}^{n+\frac12} + \frac{r_{z2}}{\Delta z}\Big(E_{y,i,j+\tfrac12,k+1}^{n+\frac12}-E_{y,i,j+\tfrac12,k}^{n+\frac12}\Big) - \frac{r_{z2}}{\Delta y}\Big(E_{z,i,j+1,k+\tfrac12}^{n+1}-E_{z,i,j,k+\tfrac12}^{n+1}\Big).
    \label{eq:adi-second-hx}
\end{equation}
where
\begin{align}
    \alpha_{z2} & =\frac{D_{b,i,j-\tfrac12,k+\tfrac12}}{\Delta y^2},       \\
    \beta_{z2}  & =\frac{1}{C_{b,i,j,k+\tfrac12}}+\alpha_{z2}+\gamma_{z2}, \\
    \gamma_{z2} & =\frac{D_{b,i,j+\tfrac12,k+\tfrac12}}{\Delta y^2},       \\
    p_{z2}      & =\frac{C_{a,i,j,k+\tfrac12}}{C_{b,i,j,k+\tfrac12}},      \\
    q_{z2}      & =D_{b,i,j-\tfrac12,k+\tfrac12},                          \\
    r_{z2}      & =D_{b,i,j+\tfrac12,k+\tfrac12}.
\end{align}

In the AMReX Yee grid, with vacuum media, the field updates are implemented as
\begin{subequations}\label{eq:adi-second-half-amrex}
    \begin{align}
        \begin{split}
            - E_x^{n+1}(i, j, k - 1) + \left(2+\frac{4\Delta z^2}{c^2\Delta t^2}\right) E_x^{n+1}(i, j, k) & - E_x^{n+1}(i, j, k + 1) = \frac{4\Delta z^2}{c^2\Delta t^2} E_x^{n+\frac12}(i, j, k)                                                                                                      \\
                                                                                                           & + \frac{2\Delta z^2}{\Delta t}\left[\frac{B_z^{n+\frac12}(i, j, k)-B_z^{n+\frac12}(i, j - 1, k)}{\Delta y} - \frac{B_y^{n+\frac12}(i, j, k)-B_y^{n+\frac12}(i, j, k - 1)}{\Delta z}\right] \\
                                                                                                           & + \Delta z\left[\frac{E_z^{n+\frac12}(i + 1, j, k - 1)-E_z^{n+\frac12}(i, j, k - 1)}{\Delta x} - \frac{E_z^{n+\frac12}(i + 1, j, k)-E_z^{n+\frac12}(i, j, k)}{\Delta x}\right],            \\
        \end{split}
    \end{align}
    \begin{align}
        \begin{split}
            - E_y^{n+1}(i - 1, j, k) + \left(2+\frac{4\Delta x^2}{c^2\Delta t^2}\right) E_y^{n+1}(i, j, k) & - E_y^{n+1}(i + 1, j, k) = \frac{4\Delta x^2}{c^2\Delta t^2} E_y^{n+\frac12}(i, j, k)                                                                                                      \\
                                                                                                           & + \frac{2\Delta x^2}{\Delta t}\left[\frac{B_x^{n+\frac12}(i, j, k)-B_x^{n+\frac12}(i, j, k - 1)}{\Delta z} - \frac{B_z^{n+\frac12}(i, j, k)-B_z^{n+\frac12}(i - 1, j, k)}{\Delta x}\right] \\
                                                                                                           & + \Delta x\left[\frac{E_x^{n+\frac12}(i - 1, j + 1, k)-E_x^{n+\frac12}(i - 1, j, k)}{\Delta y} - \frac{E_x^{n+\frac12}(i, j + 1, k)-E_x^{n+\frac12}(i, j, k)}{\Delta y}\right],            \\
        \end{split}
    \end{align}
    \begin{align}
        \begin{split}
            - E_z^{n+1}(i, j - 1, k) + \left(2+\frac{4\Delta y^2}{c^2\Delta t^2}\right) E_z^{n+1}(i, j, k) & - E_z^{n+1}(i, j + 1, k) = \frac{4\Delta y^2}{c^2\Delta t^2} E_z^{n+\frac12}(i, j, k)                                                                                                      \\
                                                                                                           & + \frac{2\Delta y^2}{\Delta t}\left[\frac{B_y^{n+\frac12}(i, j, k)-B_y^{n+\frac12}(i - 1, j, k)}{\Delta x} - \frac{B_x^{n+\frac12}(i, j, k)-B_x^{n+\frac12}(i, j - 1, k)}{\Delta y}\right] \\
                                                                                                           & + \Delta y\left[\frac{E_y^{n+\frac12}(i, j - 1, k + 1)-E_y^{n+\frac12}(i, j - 1, k)}{\Delta z} - \frac{E_y^{n+\frac12}(i, j, k + 1)-E_y^{n+\frac12}(i, j, k)}{\Delta z}\right],            \\
        \end{split}
    \end{align}
    \begin{align}
        B_x^{n+1}(i, j, k)
         & = B_x^{n+\frac12}(i, j, k) + \frac{\Delta t}{2}\left[\frac{E_y^{n+\frac12}(i, j, k + 1)-E_y^{n+\frac12}(i, j, k)}{\Delta z} - \frac{E_z^{n+1}(i, j + 1, k)-E_z^{n+1}(i, j, k)}{\Delta y}\right], \\
        B_y^{n+1}(i, j, k)
         & = B_y^{n+\frac12}(i, j, k) + \frac{\Delta t}{2}\left[\frac{E_z^{n+\frac12}(i + 1, j, k)-E_z^{n+\frac12}(i, j, k)}{\Delta x} - \frac{E_x^{n+1}(i, j, k + 1)-E_x^{n+1}(i, j, k)}{\Delta z}\right], \\
        B_z^{n+1}(i, j, k)
         & = B_z^{n+\frac12}(i, j, k) + \frac{\Delta t}{2}\left[\frac{E_x^{n+\frac12}(i, j + 1, k)-E_x^{n+\frac12}(i, j, k)}{\Delta y} - \frac{E_y^{n+1}(i + 1, j, k)-E_y^{n+1}(i, j, k)}{\Delta x}\right].
    \end{align}
\end{subequations}


\section{Solution of Linear Systems}

See \href{https://en.wikipedia.org/wiki/Tridiagonal_matrix_algorithm}{Tridiagonal matrix algorithm} for more details.

\section{Perfect Electric Conductor Boundary Conditions}
\label{sec:pec-boundary-conditions}

For a perfect electric conductor (PEC) boundary with outward unit normal $\hat{\mathbf{n}}$, the electric field satisfies
\begin{equation}
    \hat{\mathbf{n}} \times \mathbf{E} = 0.
\end{equation}
That is, the tangential electric field vanishes on the conducting surface. If the field inside the conductor is not evolved and the initial data are consistent with the boundary, then Faraday's law also implies that the normal magnetic flux remains zero on the wall,
\begin{equation}
    \hat{\mathbf{n}} \cdot \mathbf{B} = 0.
\end{equation}

Now consider a rectangular domain with PEC boundaries on the two faces normal to $x$, namely the planes $x=x_{\mathrm{lo}}$ and $x=x_{\mathrm{hi}}$, while the four faces normal to $y$ and $z$ remain periodic. On the AMReX Yee grid, $E_y$ and $E_z$ are nodal in $x$ and therefore live directly on the $x$-boundary planes, whereas $E_x$ is centered at $i+\tfrac12$ in the $x$ direction and is not stored on the boundary itself. If the domain has $N_x$ cells in $x$, the discrete PEC conditions are therefore
\begin{equation}\label{eq:pec-x-walls}
    E_y(0,j,k) = E_y(N_x,j,k) = 0, \qquad E_z(0,j,k) = E_z(N_x,j,k) = 0.
\end{equation}
No value is imposed directly on $E_x$ at the wall, since it is the normal component.

In the Yee update, the implementation is to overwrite the tangential electric fields on the two PEC planes after each electric-field stage,
\begin{equation}
    E_y^m(0,j,k) = E_y^m(N_x,j,k) = 0, \qquad E_z^m(0,j,k) = E_z^m(N_x,j,k) = 0,
\end{equation}
where $m=n+\frac12$ or $n+1$ depending on the stage. The neighboring magnetic updates then use differences against these zero boundary values. For example, the $x$-derivative appearing in the $B_y$ update is approximated by
\begin{equation}
    \left.\frac{\partial E_z}{\partial x}\right|_{i+\frac12,j,k} \approx \frac{E_z(i+1,j,k)-E_z(i,j,k)}{\Delta x}.
\end{equation}
At the first interior face next to $x=x_{\mathrm{lo}}$,
\begin{equation}
    \left.\frac{\partial E_z}{\partial x}\right|_{\frac12,j,k} \approx \frac{E_z(1,j,k)-E_z(0,j,k)}{\Delta x} = \frac{E_z(1,j,k)}{\Delta x},
\end{equation}
while at the last interior face next to $x=x_{\mathrm{hi}}$,
\begin{equation}
    \left.\frac{\partial E_z}{\partial x}\right|_{N_x-\frac12,j,k} \approx \frac{E_z(N_x,j,k)-E_z(N_x-1,j,k)}{\Delta x} = -\frac{E_z(N_x-1,j,k)}{\Delta x}.
\end{equation}
Likewise, the update for the normal magnetic component $B_x$ on the PEC planes involves only tangential derivatives of $E_y$ and $E_z$, so if $B_x$ is initialized consistently on the wall then it remains zero there.

In the ADI scheme, the same PEC condition enters as a Dirichlet condition whenever a tangential electric component is solved implicitly along $x$. Because the $y$ and $z$ directions remain periodic, the implicit solves along those directions are unchanged and stay cyclic. The only modified line solves are the $x$-directed ones:
\begin{itemize}
    \item the $E_z$ solve in the first half-step, equation \eqref{eq:first-e-ez-adi-tridiagonal};
    \item the $E_y$ solve in the second half-step, equation \eqref{eq:second-e-ey-adi-tridiagonal}.
\end{itemize}

For the first half-step, the $E_z$ solve uses
\begin{equation}
    E_{z,0,j,k+\tfrac12}^{n+\frac12} = 0, \qquad E_{z,N_x,j,k+\tfrac12}^{n+\frac12} = 0.
\end{equation}
Hence the first and last interior tridiagonal equations are
\begin{equation}
    \beta_{z1} E_{z,1,j,k+\tfrac12}^{n+\frac12} - \gamma_{z1} E_{z,2,j,k+\tfrac12}^{n+\frac12} = \mathrm{RHS},
\end{equation}
and
\begin{equation}
    -\alpha_{z1} E_{z,N_x-2,j,k+\tfrac12}^{n+\frac12} + \beta_{z1} E_{z,N_x-1,j,k+\tfrac12}^{n+\frac12} = \mathrm{RHS},
\end{equation}
respectively. The boundary unknowns do not participate in the solve; they are prescribed by the PEC condition.

Similarly, in the second half-step the $E_y$ solve uses
\begin{equation}
    E_{y,0,j+\tfrac12,k}^{n+1} = 0, \qquad E_{y,N_x,j+\tfrac12,k}^{n+1} = 0.
\end{equation}
Hence the first and last interior tridiagonal equations are
\begin{equation}
    \beta_{y2} E_{y,1,j+\tfrac12,k}^{n+1} - \gamma_{y2} E_{y,2,j+\tfrac12,k}^{n+1} = \mathrm{RHS},
\end{equation}
and
\begin{equation}
    -\alpha_{y2} E_{y,N_x-2,j+\tfrac12,k}^{n+1} + \beta_{y2} E_{y,N_x-1,j+\tfrac12,k}^{n+1} = \mathrm{RHS},
\end{equation}
respectively.

Note that we do not need to explicitly impose $\hat{\mathbf{n}} \cdot \mathbf{B} = 0$ separately. For an $x$-normal wall, the magnetic update on the boundary plane $i=i_b$, with $i_b \in \{0,N_x\}$, is
\begin{equation}
    B_x^{n+1}(i_b,j,k)
    = B_x^n(i_b,j,k)
    + \frac{\Delta t}{2}\left[
    \frac{E_y^{n+\frac12}(i_b,j,k+1)-E_y^{n+\frac12}(i_b,j,k)}{\Delta z}
    - \frac{E_z^{n+1}(i_b,j+1,k)-E_z^{n+1}(i_b,j,k)}{\Delta y}
    \right].
\end{equation}
Since the PEC condition gives
\begin{equation}
    E_y(i_b,j,k) = 0, \qquad E_z(i_b,j,k) = 0
\end{equation}
for all $j$ and $k$, it follows that
\begin{equation}
    B_x^{n+1}(i_b,j,k) = B_x^n(i_b,j,k).
\end{equation}
Therefore, if
\begin{equation}
    B_x^0(i_b,j,k) = 0,
\end{equation}
then
\begin{equation}
    B_x^n(i_b,j,k) = 0
\end{equation}
is already satisfied for all $n$.

\section{Dispersion Analysis}

We analyze the homogeneous-vacuum ADI scheme for a wave propagating in the $x$-direction with
\begin{equation}
    E_x=E_y=0,\qquad H_x=H_z=0,
\end{equation}
so that the only nontrivial fields are $E_z$ and $H_y$. The other axis-aligned polarizations follow by a cyclic permutation of coordinates.

For simplicity, we consider the case of vacuum media. Set $\sigma=0$, $\varepsilon=\varepsilon_0$, $\mu=\mu_0$, and
\begin{equation}
    c=\frac{1}{\sqrt{\varepsilon_0\mu_0}}.
\end{equation}
The half-step coefficients \eqref{eq:adi-half-coefficients} are then
\begin{equation}
    C_a=1,\qquad
    C_b=\frac{\Delta t}{2\varepsilon_0},\qquad
    D_b=\frac{\Delta t}{2\mu_0}.
    \label{eq:adi-dispersion-vacuum-coefficients}
\end{equation}
Define the coefficient $\gamma$ by
\begin{equation}
    \gamma:=C_bD_b=\frac{(c\Delta t)^2}{4}.
\end{equation}

Define the two staggered first-difference operators by
\begin{subequations}
    \begin{align}
        \left(\delta_x^{E}E_z\right)_{i+\tfrac12,j,k+\tfrac12}
         & :=\frac{E_{z,i+1,j,k+\tfrac12}-E_{z,i,j,k+\tfrac12}}{\Delta x},
        \label{eq:adi-delta-x-E}                                                           \\
        \left(\delta_x^{H}H_y\right)_{i,j,k+\tfrac12}
         & :=\frac{H_{y,i+\tfrac12,j,k+\tfrac12}-H_{y,i-\tfrac12,j,k+\tfrac12}}{\Delta x}.
        \label{eq:adi-delta-x-H}
    \end{align}
    \label{eq:adi-delta-x-operators}
\end{subequations}
Thus $\delta_x^{E}$ maps the $E_z$ grid to the $H_y$ grid, while $\delta_x^{H}$ maps the $H_y$ grid to the $E_z$ grid. Their composition is the centered second difference
\begin{equation}
    \left(\delta_x^2E_z\right)_{i,j,k+\tfrac12}
    :=\left(\delta_x^{H}\delta_x^{E}E_z\right)_{i,j,k+\tfrac12}
    =\frac{E_{z,i-1,j,k+\tfrac12}-2E_{z,i,j,k+\tfrac12}
    +E_{z,i+1,j,k+\tfrac12}}{\Delta x^2}.
    \label{eq:adi-delta-x-squared}
\end{equation}

Restricting \eqref{eq:first-e-ez-adi-tridiagonal}--\eqref{eq:adi-first-hy} to this polarization and multiplying the electric equation by $C_b$, the first half-step is
\begin{equation}
    \left(1-\gamma\delta_x^2\right)E_z^{n+\frac12}
    =E_z^n+C_b\delta_x^{H}H_y^n,
    \label{eq:first-e-ez-adi-tridiagonal-dispersion}
\end{equation}
\begin{equation}
    H_y^{n+\frac12}=H_y^n+D_b\delta_x^{E}E_z^{n+\frac12},
    \label{eq:adi-first-hy-dispersion}
\end{equation}
where $E_z$ is implicit along $x$. In the second half-step, $E_z$ is formally implicit along $y$. Because this mode has no $y$-dependence, that second-difference term vanishes and the update reduces to
\begin{equation}
    E_z^{n+1}=E_z^{n+\frac12}+C_b\delta_x^{H}H_y^{n+\frac12},
    \label{eq:second-e-ez-adi-tridiagonal-dispersion}
\end{equation}
\begin{equation}
    H_y^{n+1}=H_y^{n+\frac12}+D_b\delta_x^{E}E_z^{n+\frac12}.
    \label{eq:adi-second-hy-dispersion}
\end{equation}
Combining \eqref{eq:adi-first-hy-dispersion} and \eqref{eq:adi-second-hy-dispersion}, we have
\begin{equation}
    H_y^{n+1}=H_y^n+2D_b\delta_x^{E}E_z^{n+\frac12}
    \label{eq:adi-dispersion-full-hy}
\end{equation}

At any time level $m$, write a spatial Fourier mode on the staggered Yee grid as
\begin{align}
    E_{z,i,j,k+\tfrac12}^{m}          & =\widehat E_z^{\,m}e^{\mathrm{i}k_xi\Delta x},            \\
    H_{y,i+\tfrac12,j,k+\tfrac12}^{m} & =\widehat H_y^{\,m}e^{\mathrm{i}k_x(i+\tfrac12)\Delta x}.
    \label{eq:adi-dispersion-spatial-mode}
\end{align}
Plugging into \eqref{eq:adi-delta-x-operators} and \eqref{eq:adi-delta-x-squared} gives
\begin{subequations}
    \begin{align}
        \left(\delta_x^{E}E_z^m\right)_{i+\tfrac12,j,k+\tfrac12} & = \mathrm{i}K_x\widehat E_z^{\,m}e^{\mathrm{i}k_x(i+\tfrac12)\Delta x}, \\
        \left(\delta_x^{H}H_y^m\right)_{i,j,k+\tfrac12}          & = \mathrm{i}K_x\widehat H_y^{\,m}e^{\mathrm{i}k_xi\Delta x},            \\
        \left(\delta_x^2E_z^m\right)_{i,j,k+\tfrac12}            & =-K_x^2\widehat E_z^{\,m}e^{\mathrm{i}k_xi\Delta x}.
    \end{align}
    \label{eq:adi-delta-x-fourier}
\end{subequations}
where
\begin{equation}
    K_x:=\frac{2}{\Delta x}\sin\!\left(\frac{k_x\Delta x}{2}\right).
    \label{eq:adi-Kx}
\end{equation}

For brevity, define the derivative symbol and implicit denominator
\begin{equation}
    s:=\mathrm{i}K_x,\qquad D:=1+\gamma K_x^2.
    \label{eq:adi-dispersion-s-D}
\end{equation}
Substituting into \eqref{eq:first-e-ez-adi-tridiagonal-dispersion} and using \eqref{eq:adi-delta-x-fourier}, the first-half electric update becomes
\begin{equation}
    D\widehat E_z^{\,n+\frac12} = \widehat E_z^{\,n}+C_bs\widehat H_y^{\,n}.
    \label{eq:adi-dispersion-half1-E}
\end{equation}
Similarly for \eqref{eq:adi-first-hy-dispersion}, we have
\begin{equation}
    \widehat H_y^{\,n+\frac12} = \widehat H_y^{\,n}+D_bs\widehat E_z^{\,n+\frac12}.
    \label{eq:adi-dispersion-half1-H}
\end{equation}
Substituting into \eqref{eq:second-e-ez-adi-tridiagonal-dispersion} and using \eqref{eq:adi-delta-x-fourier}, the second-half electric update becomes
\begin{equation}
    \widehat E_z^{\,n+1} = \widehat E_z^{\,n+\frac12}+C_bs\widehat H_y^{\,n+\frac12},
    \label{eq:adi-dispersion-half2-E}
\end{equation}
Similarly for \eqref{eq:adi-second-hy-dispersion}, we have
\begin{equation}
    \widehat H_y^{\,n+1} = \widehat H_y^{\,n}+2D_bs\widehat E_z^{\,n+\frac12}
    \label{eq:adi-dispersion-full-H}
\end{equation}

Solving \eqref{eq:adi-dispersion-half1-E} for the intermediate electric amplitude gives
\begin{equation}
    \widehat E_z^{\,n+\frac12}=\frac{1}{D}\widehat E_z^{\,n}+\frac{C_bs}{D}\widehat H_y^{\,n}.
\end{equation}
Substitution into \eqref{eq:adi-dispersion-half2-E} and \eqref{eq:adi-dispersion-full-H}, using $C_bD_b=\gamma$ and $s^2=-K_x^2$, gives the one-step map
\begin{equation}
    \begin{pmatrix} \widehat E_z \\ \widehat H_y \end{pmatrix}^{n+1}=G\begin{pmatrix} \widehat E_z \\ \widehat H_y \end{pmatrix}^{n},\qquad G=\frac{1}{D}\begin{pmatrix} 1-\gamma K_x^2 & 2C_b s \\ 2D_b s & 1-\gamma K_x^2 \end{pmatrix}.
    \label{eq:adi-amplification-matrix}
\end{equation}

A time-harmonic mode satisfies
\begin{equation}
    \begin{pmatrix}\widehat E_z\\\widehat H_y\end{pmatrix}^{n+1}=\lambda\begin{pmatrix}\widehat E_z\\\widehat H_y\end{pmatrix}^{n},\qquad \lambda=e^{-\mathrm{i}\omega\Delta t}.
\end{equation}
The characteristic equation $\det(G-\lambda I)=0$ yields
\begin{equation}
    \lambda = \frac{1-\gamma K_x^2\pm 2\mathrm{i}\,K_x\sqrt{\gamma}}{1+\gamma K_x^2}.
\end{equation}
Introduce the CFL number $S:=c\Delta t/\Delta x$ and
\begin{equation}
    \nu := S\sin\!\left(\frac{k_x\Delta x}{2}\right) = K_x\sqrt{\gamma}.
\end{equation}
Then $\gamma K_x^2=\nu^2$, and the eigenvalues simplify to
\begin{equation}
    \lambda = \frac{1\pm\mathrm{i}\nu}{1\mp\mathrm{i}\nu}.
\end{equation}
The forward-propagating branch ($\omega>0$) is
\begin{equation}
    \lambda = \frac{1-\mathrm{i}\nu}{1+\mathrm{i}\nu},
\end{equation}
for which $\lvert\lambda\rvert=1$ indicating unconditional stability of the ADI scheme and there is no numerical dispersion. The phase angle is
\begin{equation}
    \arg(\lambda)=-2\arctan(\nu),
\end{equation}
so the ADI numerical angular frequency is
\begin{equation}
    \omega_{\mathrm{ADI}}(k_x) = \frac{2}{\Delta t}\,\arctan\!\left(S\sin\frac{k_x\Delta x}{2}\right).
    \label{eq:adi-dispersion-relation}
\end{equation}
The corresponding numerical phase velocity is
\begin{equation}
    v_{p,\mathrm{ADI}}(k_x)=\frac{\omega_{\mathrm{ADI}}(k_x)}{k_x},
\end{equation}
which generally differs from $c$ and therefore exhibits numerical dispersion.

\subsection{Numerical experiments}

\paragraph{Initial-value problem verification}
The numerical dispersion relation was verified with a PEC standing wave and a periodic traveling wave. Both cases used the macroscopic vacuum ADI solver with the target resolution
\begin{equation}
    k\Delta z \simeq 2.15\times10^{-3}
\end{equation}
and CFL numbers $S=c\Delta t/\Delta z=64$, $128$, $256$, and $512$. In both cases $N_x=N_y=8$. Each simulation was evolved for sixteen periods of the continuum frequency $f_0=ck/(2\pi)$. The surface integral $\int E_y\,\mathrm{d}x\,\mathrm{d}y$ was sampled on a surface normal to $z$ at $z=L_z/4$, approximately sixteen times per period. The dominant frequency was extracted from the last half of the signal using a Hann window, eightfold zero padding, and normalization by the number of samples.
\begin{figure}[htpb]
    \centering
    \begin{subfigure}[t]{0.3\textwidth}
        \centering
        \includegraphics[width=\linewidth]{png_out/pec_standing_dispersion_vs_cfl.png}
        \caption{Normalized frequency versus CFL number.}
        \label{fig:adi-phase-velocity-standing-dispersion}
    \end{subfigure}
    \hfill
    \begin{subfigure}[t]{0.65\textwidth}
        \centering
        \includegraphics[width=\linewidth]{png_out/pec_standing_resonance_vs_cfl.png}
        \caption{Time histories and corresponding frequency spectra.}
        \label{fig:adi-phase-velocity-standing-spectra}
    \end{subfigure}
    \caption{Numerical verification of the ADI dispersion relation using PEC standing and periodic traveling waves.}
    \label{fig:adi-phase-velocity-standing}
\end{figure}

Both tests used $L_z=8\,\mu\mathrm{m}$ and $N_z=2928$, corresponding to one wavelength across the domain:
\begin{equation}
    k=\frac{2\pi}{L_z},\qquad
    k\Delta z=\frac{2\pi}{N_z}\simeq2.146\times10^{-3},\qquad
    f_0\simeq3.75\times10^{13}\ \mathrm{Hz}.
\end{equation}
For the PEC case, perfect-electric-conductor boundaries were imposed at $z=0$ and $z=L_z$, with periodic boundaries in $x$ and $y$, and the standing wave was initialized as
\begin{equation}
    E_y=E_0\sin(2\pi z/L_z),\qquad \mathbf{B}=0.
\end{equation}
For the periodic case, periodic boundaries were imposed in all directions and the traveling wave was initialized as
\begin{equation}
    E_y=E_0\sin(2\pi z/L_z),\qquad B_x=-E_y/c.
\end{equation}

The measured FFT peak was divided by $f_0$. For this fixed spatial mode the normalized frequency is also the normalized phase velocity, $f/f_0=v_p/c$, and the analytic ADI prediction is
\begin{equation}
    \frac{f_{\mathrm{ADI}}}{f_0}
    =\frac{2\arctan\!\left[S\sin(k\Delta z/2)\right]}{S\,k\Delta z}.
\end{equation}


Figure~\ref{fig:adi-phase-velocity-standing} compares the measured frequencies with the analytic dispersion relation and shows the corresponding time histories and spectra. Both the periodic traveling-wave case and the PEC standing-wave case show good agreement with the analytical frequency shift, verifying the ADI dispersion relation. The agreement of both initial-value problems also indicates that the PEC boundary does not introduce observable additional error.

\paragraph{Soft excitation verification}

The dispersion relation was also tested by exciting the PEC cavity from initially zero fields. The geometry, material parameters, resolution, CFL numbers, run duration, and boundary conditions were identical to those of the PEC standing-wave test above. Only $E_y$ was excited, using the soft source
\begin{equation}
    S_{E_y}(z,t)
    =E_0\frac{\Delta t}{T_P}
    \sin\!\left(\frac{2\pi z}{L_z}\right)
    \exp\!\left[-\frac{(t-3T_P)^2}{2T_P^2}\right]
    \sin(2\pi f_0t),
\end{equation}
where $E_0=1$, $f_0=c/L_z$, and $T_P=1/f_0$. Thus the source has the spatial profile of the one-wavelength cavity mode and a Gaussian envelope centered at $3T_P$ with a width of one period. The factor $\Delta t/T_P$ keeps the total excitation approximately independent of the CFL number. No magnetic excitation was applied.
\begin{figure}[htpb]
    \centering
    \begin{subfigure}[t]{0.3\textwidth}
        \centering
        \includegraphics[width=\linewidth]{png_out/pec_soft_dispersion_vs_cfl.png}
        \caption{Normalized frequency versus CFL number.}
        \label{fig:adi-phase-velocity-soft-dispersion}
    \end{subfigure}
    \hfill
    \begin{subfigure}[t]{0.65\textwidth}
        \centering
        \includegraphics[width=\linewidth]{png_out/pec_soft_resonance_vs_cfl.png}
        \caption{Time histories and corresponding frequency spectra.}
        \label{fig:adi-phase-velocity-soft-spectra}
    \end{subfigure}
    \caption{Numerical verification of the ADI dispersion relation using PEC soft excitation and periodic traveling waves.}
    \label{fig:adi-phase-velocity-soft}
\end{figure}

The electric field was sampled at the mode antinode $z=L_z/4$ at approximately sixteen samples per period. As in the initial-value tests, the dominant frequency was extracted from the last half of the sixteen-period signal using a Hann-windowed, zero-padded Fourier transform. The measured frequencies were compared with \eqref{eq:adi-dispersion-relation} for $S=64$, $128$, $256$, and $512$. A periodic traveling-wave initial-value problem on the same grid provided a reference.


\paragraph{The origin of the amplitude variation}
After the finite-duration source has vanished, the cavity oscillates at its intrinsic numerical frequency $\omega_{\mathrm{ADI}}(S)$. The experiments show that $\omega_{\mathrm{ADI}}(S)$ is independent of the excitation frequency $\omega_d$. The residual amplitude, however, is determined by the spectral overlap of the source pulse with this cavity mode. For the Gaussian-modulated sinusoidal source used here, the positive-frequency contribution is approximately
\begin{equation}
    A(S)\propto
    \exp\!\left[
        -\frac{T_P^2}{2}
        \bigl(\omega_{\mathrm{ADI}}(S)-\omega_d\bigr)^2
        \right].
\end{equation}

Besides this detuning effect, the finite temporal sampling of the soft source changes the amount of drive injected into the cavity. The carrier is sampled once per time step with spacing $\Delta t=S\Delta z/c$. Representing each sampled value over one time-step interval gives the discrete-injection factor
\begin{equation}
    \operatorname{sinc}\!\left(\frac{\omega_d\Delta t}{2}\right)
    =
    \frac{\sin(\omega_d\Delta t/2)}{\omega_d\Delta t/2}.
\end{equation}
This factor is separate from the ADI dispersion relation: it describes the loss of effective drive amplitude caused by resolving the carrier with a finite number of time steps. Combining the two effects gives the amplitude model
\begin{equation}
    A(S)\propto
    \underbrace{
        \exp\!\left[
            -\frac{T_P^2}{2}
            \bigl(\omega_{\mathrm{ADI}}(S)-\omega_d\bigr)^2
            \right]
    }_{\text{drive--mode detuning}}
    \underbrace{
        \operatorname{sinc}\!\left(\frac{\omega_d\Delta t}{2}\right)
    }_{\text{discrete drive injection}}.
\end{equation}
For the present grid, increasing the CFL number from $64$ to $512$ reduces the number of time steps per drive period from approximately $46$ to $6$. At $S=512$, $\omega_d\Delta t/2=\pi S/N_z\simeq0.55$ and the injection factor is approximately $0.95$. The detuning factor alone predicts an amplitude ratio of approximately $0.865$, whereas inclusion of the discrete-injection factor gives $0.822$, close to the measured value $0.827$. This correction applies only to the externally driven case and is absent from the standing-wave initial-value problem.

Because $\omega_{\mathrm{ADI}}(S)$ decreases as the CFL number increases, the response grows when this shift moves the numerical cavity frequency toward $\omega_d$ and decreases when it moves the cavity frequency away from $\omega_d$. If the numerical frequency crosses the excitation frequency, the amplitude reaches a maximum near the crossing and the trend reverses. The observed CFL dependence therefore reflects drive--mode detuning, rather than numerical damping of the ADI mode; temporal sampling of the source may introduce an additional amplitude error at the largest CFL numbers.

Figure~\ref{fig:adi-amplitude-soft-vs-standing-ref} compares the measured soft-excitation response with the detuning and discrete-injection models. The soft-excitation amplitude curve exactly matches our detuning and discrete-injection model.
\begin{figure}[htpb]
    \centering
    \includegraphics[width=\linewidth]{png_out/pec_soft_vs_standing_cfl_ref.png}
    \caption{Amplitude and frequency of the PEC soft-excitation and standing-wave responses versus CFL number. The amplitude reference curves include the detuning-only model and the combined detuning and discrete-injection model.}
    \label{fig:adi-amplitude-soft-vs-standing-ref}
\end{figure}


\section{External Excitation}

External field excitations are prescribed by user-defined grid functions of space and time, together with a spatial flag that selects the excitation type at each Yee edge (or face). Following the Artemis convention,
\begin{itemize}
    \item flag $=0$: no excitation,
    \item flag $=1$: hard source (Dirichlet), field equals the prescribed value,
    \item flag $=2$: soft source, the prescribed value is \emph{added} to the field update.
\end{itemize}
In the macroscopic ADI push these two families are treated differently: soft electric sources are built into the implicit electric systems, while magnetic sources are applied after each magnetic half-step. Hard electric sources are not supported in the ADI electric solve.

\subsection{Soft Electric Sources in the ADI Electric Update}

Write $S_{E_x,i+\tfrac12,j,k}^{n+\frac12}$ for the soft-source value assigned to the $E_x$ edge at $\bigl(i+\tfrac12,j,k\bigr)$ and the arrival time level $n+\tfrac12$ (and likewise for $E_y$ and $E_z$ at their staggered locations). Over a full step the intended soft contribution at that edge is split equally between the two ADI half-steps and evaluated at the resulting field levels $n+\tfrac12$ and $n+1$. The first-half electric update \eqref{eq:first-e-ex-adi} is then modified to
\begin{equation}\label{eq:adi-ex-soft-unnormalized}
    E_{x,i+\tfrac12,j,k}^{n+\frac12}
    = C_{a,i+\tfrac12,j,k}\,E_{x,i+\tfrac12,j,k}^{n}
    + C_{b,i+\tfrac12,j,k}\,[H_1]
    + \frac12\,S_{E_x,i+\tfrac12,j,k}^{n+\frac12},
\end{equation}
on edges where the $E_x$ flag equals $2$, and is unchanged elsewhere. The second half-step \eqref{eq:second-e-ex-adi} is analogous,
\begin{equation}
    E_{x,i+\tfrac12,j,k}^{n+1}
    = C_{a,i+\tfrac12,j,k}\,E_{x,i+\tfrac12,j,k}^{n+\frac12}
    + C_{b,i+\tfrac12,j,k}\,[H_4]
    + \frac12\,S_{E_x,i+\tfrac12,j,k}^{n+1}.
\end{equation}
The same construction applies componentwise to $E_y$ and $E_z$. The magnetic updates themselves are unchanged by the electric source; only the electric right-hand side is affected.

\subsection{Effect on the ADI Tridiagonal System}

Dividing \eqref{eq:adi-ex-soft-unnormalized} by $C_{b,i+\tfrac12,j,k}$ and substituting the magnetic update as in Section~3 yields the same left-hand side as \eqref{eq:first-e-ex-adi-tridiagonal}. On soft-source edges the only modification is an additional known term on the right-hand side:
\begin{align}
    \begin{split}
        -\alpha_{x1}E_{x,i+\tfrac12,j-1,k}^{n+\frac12}
        + \beta_{x1}E_{x,i+\tfrac12,j,k}^{n+\frac12}
        - \gamma_{x1}E_{x,i+\tfrac12,j+1,k}^{n+\frac12}
         & = p_{x1}E_{x,i+\tfrac12,j,k}^{n}                                                                       \\
         & \quad +\frac{1}{\Delta y}\Big(H_{z,i+\tfrac12,j+\tfrac12,k}^{n}-H_{z,i+\tfrac12,j-\tfrac12,k}^{n}\Big) \\
         & \quad -\frac{1}{\Delta z}\Big(H_{y,i+\tfrac12,j,k+\tfrac12}^{n}-H_{y,i+\tfrac12,j,k-\tfrac12}^{n}\Big) \\
         & \quad +\frac{q_{x1}}{\Delta x\Delta y}\Big(E_{y,i+1,j-\tfrac12,k}^{n}-E_{y,i,j-\tfrac12,k}^{n}\Big)    \\
         & \quad -\frac{r_{x1}}{\Delta x\Delta y}\Big(E_{y,i+1,j+\tfrac12,k}^{n}-E_{y,i,j+\tfrac12,k}^{n}\Big)    \\
         & \quad +\frac{1}{2C_{b,i+\tfrac12,j,k}}\,S_{E_x,i+\tfrac12,j,k}^{n+\frac12},
    \end{split}
    \label{eq:first-e-ex-adi-tridiagonal-excitation}
\end{align}
with the same coefficients $\alpha_{x1},\beta_{x1},\gamma_{x1},p_{x1},q_{x1},r_{x1}$ as in Section~3. Away from soft-source edges the final term is omitted. In the implementation this contribution is the increment
\begin{equation}
    \texttt{rhs}_{i,j,k} \mathrel{+}= \tfrac12\,S_{E_x,i+\tfrac12,j,k}^{n+\frac12}/C_{b,i+\tfrac12,j,k}
\end{equation}
added after the curl terms have been assembled. The $E_y$ and $E_z$ first-half systems are modified in exactly the same way, and the second half-step is analogous with $S_{E_x,i+\tfrac12,j,k}^{n+1}/(2C_{b,i+\tfrac12,j,k})$ on the normalized right-hand side.

Consequently a hard electric source (flag $=1$), which would require constraining the unknown $E$ itself rather than adding to the right-hand side, is rejected in the ADI electric path.

\subsection{Magnetic Sources After Each Half-Step}

After the first ADI half-step has produced $\mathbf{E}^{n+\frac12}$ and $\mathbf{B}^{n+\frac12}$, and again after the second half-step has produced $\mathbf{E}^{n+1}$ and $\mathbf{B}^{n+1}$, an external magnetic excitation is applied directly to the magnetic field. The excitation is evaluated at the arrival level of the field being modified. Thus, after the first half-step, a soft magnetic source (flag $=2$) contributes
\begin{equation}
    B_{x,i,j+\tfrac12,k+\tfrac12}
    \;\leftarrow\;
    B_{x,i,j+\tfrac12,k+\tfrac12}
    + \tfrac12\,S_{B_x,i,j+\tfrac12,k+\tfrac12}^{n+\frac12},
\end{equation}
while a hard magnetic source (flag $=1$) overwrites
\begin{equation}
    B_{x,i,j+\tfrac12,k+\tfrac12}
    \;\leftarrow\;
    S_{B_x,i,j+\tfrac12,k+\tfrac12}^{n+\frac12}.
\end{equation}
After the second half-step the corresponding source values are $S_{B_x}^{n+1}$: a soft source adds $\tfrac12 S_{B_x}^{n+1}$, while a hard source overwrites the field with $S_{B_x}^{n+1}$. (The factor $\tfrac12$ again splits a full-step soft increment across the two half-steps.) The same rules apply to $B_y$ and $B_z$ at their staggered faces. Unlike the electric soft source, this magnetic update is an explicit post-processing of the field after the ADI magnetic stage; it does not enter the tridiagonal electric systems.

\subsection{Summary of the ADI Excitation Path}

Over one full step $\Delta t$, the macroscopic ADI scheme therefore proceeds as follows.
\begin{enumerate}
    \item Form and solve the first-half electric systems with soft electric sources $S_E^{n+\frac12}$ added as $\tfrac12 S_E^{n+\frac12}$ (equivalently $\tfrac12 S_E^{n+\frac12}/C_b$ on the normalized RHS).
    \item Update $\mathbf{B}$ to $n+\tfrac12$ from the ADI magnetic formulae.
    \item Apply magnetic excitation to $\mathbf{B}^{n+\frac12}$ (soft: add $\tfrac12 S_B^{n+\frac12}$; hard: set equal to $S_B^{n+\frac12}$).
    \item Form and solve the second-half electric systems with soft electric sources $S_E^{n+1}$.
    \item Update $\mathbf{B}$ to $n+1$, then apply magnetic excitation evaluated at $n+1$.
\end{enumerate}
Hard electric sources are not part of this path. Soft electric sources never overwrite $E$; they only augment the Amp\`ere right-hand side of each half-step, consistently with the half-step coefficients $C_a$ and $C_b$ of \eqref{eq:adi-half-coefficients}.

\section{Lumped Inductance}

ADI removes the Maxwell CFL restriction, so a lumped inductor must not reintroduce a local timestep limit from the $LC$ scale $\omega_L\sim\sqrt{\kappa/\varepsilon}$. This section embeds an ideal inductor in the ADI electric update by a centered trapezoidal coupling: it is nondissipative on the isolated field--inductor subsystem and places no restriction on $\omega_L\Delta t$.

\subsection{Continuum model}

A lumped inductor enters Amp\`ere's law as an additional current density $\mathbf{J}_\mathrm{L}$,
\begin{subequations}\label{eq:maxwell-inductor}
    \begin{align}
        \nabla \times \mathbf{H} & = \varepsilon \frac{\partial \mathbf{E}}{\partial t} + \sigma \mathbf{E} + \mathbf{J}_\mathrm{L}
        \label{eq:maxwell-ampere-inductor}                                                                                          \\
        \nabla \times \mathbf{E} & = -\mu \frac{\partial \mathbf{H}}{\partial t}.
        \label{eq:maxwell-faraday-inductor}
    \end{align}
\end{subequations}
On each Yee edge that hosts an inductor, the constitutive relation $V=L\,dI/dt$ is written in field form. For an $x$-directed edge of length $\Delta x$ and face area $\Delta y\,\Delta z$, with total inductance $L_x$ assigned to that edge,
\begin{equation}
    V = E_x\Delta x,
    \qquad
    I = J_{\mathrm{L},x}\,\Delta y\,\Delta z,
\end{equation}
so
\begin{equation}\label{eq:inductor-constitutive}
    E_x\Delta x
    = L_x\,\Delta y\,\Delta z\,\frac{\partial J_{\mathrm{L},x}}{\partial t},
\end{equation}
or equivalently
\begin{equation}\label{eq:inductor-kappa}
    \frac{\partial J_{\mathrm{L},x}}{\partial t}
    = \kappa_x E_x,
    \qquad
    \kappa_x = \frac{\Delta x}{\Delta y\,\Delta z\,L_x}.
\end{equation}
The $y$ and $z$ components are analogous, with
\begin{equation}
    \kappa_y = \frac{\Delta y}{\Delta x\,\Delta z\,L_y},
    \qquad
    \kappa_z = \frac{\Delta z}{\Delta x\,\Delta y\,L_z}.
\end{equation}
Away from inductor edges one sets $\kappa=0$ (equivalently $L\to\infty$), so $\mathbf{J}_\mathrm{L}$ need not be stored or evolved there; outside the inductor region it should normally be initialized to zero. Faraday's law \eqref{eq:maxwell-faraday-inductor} is unmodified, so the magnetic ADI updates of Section~3 are unchanged. Only the electric update and the inductor ODE are coupled.

\subsection{Centered trapezoidal coupling}

Over a half step of size $\Delta t/2$, discretize \eqref{eq:inductor-kappa} by the trapezoidal rule,
\begin{equation}\label{eq:inductor-trapezoidal}
    \frac{J_{\mathrm{L},x,i+\tfrac12,j,k}^{n+\frac12}-J_{\mathrm{L},x,i+\tfrac12,j,k}^{n}}{\Delta t/2}
    = \kappa_{x,i+\tfrac12,j,k}\,\frac{E_{x,i+\tfrac12,j,k}^{n+\frac12}+E_{x,i+\tfrac12,j,k}^{n}}{2}.
\end{equation}
Amp\`ere's law must use the same centered average $(J_{\mathrm{L},x}^{n+\frac12}+J_{\mathrm{L},x}^{n})/2$; using only the advanced value $J_{\mathrm{L},x}^{n+\frac12}$ is dissipative (Section~\ref{sec:inductor-energy}). With the same conductivity average as in Section~3,
\begin{equation}
    \varepsilon_{i+\tfrac12,j,k}\frac{E_{x,i+\tfrac12,j,k}^{n+\frac12}-E_{x,i+\tfrac12,j,k}^{n}}{\Delta t/2}
    + \sigma_{i+\tfrac12,j,k}\frac{E_{x,i+\tfrac12,j,k}^{n+\frac12}+E_{x,i+\tfrac12,j,k}^{n}}{2}
    + \frac{J_{\mathrm{L},x,i+\tfrac12,j,k}^{n+\frac12}+J_{\mathrm{L},x,i+\tfrac12,j,k}^{n}}{2}
        = [H_1],
    \label{eq:first-e-ex-inductor-semi}
\end{equation}
where $[H_1]$ is unchanged from Section~3. From \eqref{eq:inductor-trapezoidal},
\begin{equation}
    \frac{J_{\mathrm{L},x,i+\tfrac12,j,k}^{n+\frac12}+J_{\mathrm{L},x,i+\tfrac12,j,k}^{n}}{2}
    = J_{\mathrm{L},x,i+\tfrac12,j,k}^{n}
    + \frac{\Delta t}{8}\,\kappa_{x,i+\tfrac12,j,k}\bigl(E_{x,i+\tfrac12,j,k}^{n+\frac12}+E_{x,i+\tfrac12,j,k}^{n}\bigr).
\end{equation}
Substituting into \eqref{eq:first-e-ex-inductor-semi} and collecting unknowns gives
\begin{equation}
    \left(\frac{2\varepsilon}{\Delta t}+\frac{\sigma}{2}+\frac{\Delta t}{8}\kappa_x\right)E_{x}^{n+\frac12}
    = \left(\frac{2\varepsilon}{\Delta t}-\frac{\sigma}{2}-\frac{\Delta t}{8}\kappa_x\right)E_{x}^{n}
    + [H_1]
    - J_{\mathrm{L},x}^{n}.
\end{equation}
Define the inductor-modified half-step coefficients
\begin{equation}\label{eq:adi-half-coefficients-inductor}
    C_a^{L}
    = \frac{\dfrac{2\varepsilon}{\Delta t}-\dfrac{\sigma}{2}-\dfrac{\Delta t}{8}\kappa}
    {\dfrac{2\varepsilon}{\Delta t}+\dfrac{\sigma}{2}+\dfrac{\Delta t}{8}\kappa},
    \qquad
    C_b^{L}
    = \frac{1}
    {\dfrac{2\varepsilon}{\Delta t}+\dfrac{\sigma}{2}+\dfrac{\Delta t}{8}\kappa}.
\end{equation}
When $\kappa=0$ these reduce to \eqref{eq:adi-half-coefficients}. The electric update is then
\begin{equation}
    E_{x,i+\tfrac12,j,k}^{n+\frac12}
    = C_{a,i+\tfrac12,j,k}^{L}\,E_{x,i+\tfrac12,j,k}^{n}
    + C_{b,i+\tfrac12,j,k}^{L}\Bigl([H_1]-J_{\mathrm{L},x,i+\tfrac12,j,k}^{n}\Bigr).
    \label{eq:first-e-ex-adi-inductor-semi}
\end{equation}
Thus $\kappa$ modifies the local diagonal coefficients $C_a^{L}$ and $C_b^{L}$, while the known current $J_{\mathrm{L}}^{n}$ enters the right-hand side. After $E_x^{n+\frac12}$ is obtained, rearrange \eqref{eq:inductor-trapezoidal} to advance $J_{\mathrm{L},x}$ by
\begin{equation}
    J_{\mathrm{L},x,i+\tfrac12,j,k}^{n+\frac12}
    = J_{\mathrm{L},x,i+\tfrac12,j,k}^{n}
    + \frac{\Delta t}{4}\,\kappa_{x,i+\tfrac12,j,k}\bigl(E_{x,i+\tfrac12,j,k}^{n+\frac12}+E_{x,i+\tfrac12,j,k}^{n}\bigr).
    \label{eq:inductor-j-semi}
\end{equation}

\subsection{Local energy and why the current must be averaged}
\label{sec:inductor-energy}

On an inductor edge with $\kappa>0$, $\sigma=0$, and no curl contribution, define
\begin{equation}
    \mathcal{E}
    = \frac{\varepsilon}{2}E^2
    + \frac{1}{2\kappa}J_{\mathrm{L}}^2.
\end{equation}
The centered pair
\begin{equation}
    \varepsilon\frac{E^{n+\frac12}-E^{n}}{\Delta t/2}
    + \frac{J_{\mathrm{L}}^{n+\frac12}+J_{\mathrm{L}}^{n}}{2}
    = 0,
    \qquad
    \frac{J_{\mathrm{L}}^{n+\frac12}-J_{\mathrm{L}}^{n}}{\Delta t/2}
    = \kappa\frac{E^{n+\frac12}+E^{n}}{2},
\end{equation}
conserves this energy exactly: $\mathcal{E}^{n+\frac12}-\mathcal{E}^{n}=0$. Replacing the Amp\`ere current by the fully advanced value $J_{\mathrm{L}}^{n+\frac12}$ instead yields
\begin{equation}
    \mathcal{E}^{n+\frac12}-\mathcal{E}^{n}
    = -\frac{1}{2\kappa}\bigl(J_{\mathrm{L}}^{n+\frac12}-J_{\mathrm{L}}^{n}\bigr)^2
    \le 0,
\end{equation}
so the scheme is numerically dissipative even though the inductor ODE itself is integrated by the trapezoidal rule. That one-sided choice removes the $\omega_L\Delta t$ restriction but acts like an artificial resistance; we reject it for that reason.

\subsection{Effect on the ADI tridiagonal system}

The magnetic update for $H_z^{n+\frac12}$ is identical to Section~3. Substituting it into \eqref{eq:first-e-ex-adi-inductor-semi} yields the same stencil as \eqref{eq:first-e-ex-adi-tridiagonal}, but with $C_{a,i+\tfrac12,j,k}\to C_{a,i+\tfrac12,j,k}^{L}$ and $C_{b,i+\tfrac12,j,k}\to C_{b,i+\tfrac12,j,k}^{L}$:
\begin{align}
    \begin{split}
        -\alpha_{x1}E_{x,i+\tfrac12,j-1,k}^{n+\frac12}
        + \beta_{x1}^{L}E_{x,i+\tfrac12,j,k}^{n+\frac12}
        - \gamma_{x1}E_{x,i+\tfrac12,j+1,k}^{n+\frac12}
         & = p_{x1}^{L}E_{x,i+\tfrac12,j,k}^{n}                                                                   \\
         & \quad +\frac{1}{\Delta y}\Big(H_{z,i+\tfrac12,j+\tfrac12,k}^{n}-H_{z,i+\tfrac12,j-\tfrac12,k}^{n}\Big) \\
         & \quad -\frac{1}{\Delta z}\Big(H_{y,i+\tfrac12,j,k+\tfrac12}^{n}-H_{y,i+\tfrac12,j,k-\tfrac12}^{n}\Big) \\
         & \quad +\frac{q_{x1}}{\Delta x\Delta y}\Big(E_{y,i+1,j-\tfrac12,k}^{n}-E_{y,i,j-\tfrac12,k}^{n}\Big)    \\
         & \quad -\frac{r_{x1}}{\Delta x\Delta y}\Big(E_{y,i+1,j+\tfrac12,k}^{n}-E_{y,i,j+\tfrac12,k}^{n}\Big)    \\
         & \quad - J_{\mathrm{L},x,i+\tfrac12,j,k}^{n},
    \end{split}
    \label{eq:first-e-ex-adi-tridiagonal-inductor}
\end{align}
where
\begin{align}
    \alpha_{x1} & =\frac{D_{b,i+\tfrac12,j-\tfrac12,k}}{\Delta y^2},           \\
    \gamma_{x1} & =\frac{D_{b,i+\tfrac12,j+\tfrac12,k}}{\Delta y^2},           \\
    \beta_{x1}^{L}
                & =\frac{1}{C_{b,i+\tfrac12,j,k}^{L}}+\alpha_{x1}+\gamma_{x1}, \\
    p_{x1}^{L}
                & =\frac{C_{a,i+\tfrac12,j,k}^{L}}{C_{b,i+\tfrac12,j,k}^{L}},  \\
    q_{x1}      & =D_{b,i+\tfrac12,j-\tfrac12,k},                              \\
    r_{x1}      & =D_{b,i+\tfrac12,j+\tfrac12,k}.
\end{align}
The off-diagonal coefficients $\alpha_{x1}$ and $\gamma_{x1}$ are unchanged (they come only from the implicit curl of $H$). The diagonal $\beta_{x1}^{L}$ and the factor $p_{x1}^{L}$ feel $L_x$ through $C_{b,i+\tfrac12,j,k}^{L}$ and $C_{a,i+\tfrac12,j,k}^{L}$. The $E_y$ and $E_z$ first-half systems are modified in the same way.

The current update is
\begin{equation}
    J_{\mathrm{L},x,i+\tfrac12,j,k}^{n+\frac12}
    = J_{\mathrm{L},x,i+\tfrac12,j,k}^{n}
    + \frac{\Delta t}{4}\,\kappa_{x,i+\tfrac12,j,k}\bigl(E_{x,i+\tfrac12,j,k}^{n+\frac12}+E_{x,i+\tfrac12,j,k}^{n}\bigr).
\end{equation}

The second half-step is analogous.

\subsection{Stability}

Eliminating $J_\mathrm{L}$ locally produces an oscillator of frequency $\omega_L=\sqrt{\kappa/\varepsilon}$. On the isolated lossless field--inductor subsystem, the centered scheme of Section~\ref{sec:inductor-energy} has amplification factors of unit modulus for every $\Delta t$, so it introduces no lumped-inductor CFL. That local fact is not by itself a complete proof for the full spatial ADI splitting; unconditional stability of the coupled scheme requires that the same centered coupling be embedded consistently in an energy-stable ADI discretization.

By contrast, the explicit rule
$J_{\mathrm{L}}^{n+\frac12}=J_{\mathrm{L}}^{n}+(\Delta t/2)\kappa E^{n}$,
with the known current inserted only as a right-hand-side source, leaves $\omega_L\Delta t$ constrained and defeats the purpose of ADI at large CFL. That is the leapfrog macroscopic FDTD inductor, and it is unsuitable here.

\subsection{Present implementation}

In the current Artemis macroscopic ADI push, soft electric sources are added to the ADI right-hand side, but $\mathbf{J}_\mathrm{L}$ is not. The inductor object still advances $\mathbf{J}_\mathrm{L}$ explicitly when enabled, yet that current is unused by \texttt{MacroscopicEvolveADI}. A stable ADI inductor therefore requires: (i)~replace $C_a,C_b$ by $C_a^{L},C_b^{L}$ on inductor edges; (ii)~add $-C_b^{L}J_{\mathrm{L}}^{\mathrm{old}}$ to the field-update form, or equivalently $-J_{\mathrm{L}}^{\mathrm{old}}$ to the normalized tridiagonal right-hand side; and (iii)~update $J_{\mathrm{L}}$ with \eqref{eq:inductor-j-semi} after each electric solve.

\end{document}
